% interaction avec Linux au travers de JTAG ?

% get_user/put_user ou 
% TODO : 
%        chargement FPGA
%        communication avec GPIO en userspace, kernel space, et par devicetree overlay avec
%            l'appel automatique au module

% also at https://github.com/RedPitaya/RedPitaya/blob/master/fpga/fpga.rst

% make dtbs dans le noyau linux pour refaire le devicetree

\documentclass{article}
\usepackage{ucs,multicol}
\usepackage{url}
\usepackage{xcolor}
\usepackage{lastpage,graphicx}
\usepackage{geometry} % :,fourier,hyperref}
\usepackage{listings,framed}
\geometry{verbose,a4paper,tmargin=2.5cm,bmargin=2cm,lmargin=2cm,rmargin=1cm}
\usepackage{fancyhdr}
\pagestyle{fancy}
\setlength\parskip{\medskipamount}
\setlength\parindent{0pt}

\usepackage{listings}
\usepackage{color,picinpar}

\definecolor{grey}{rgb}{0.95,0.95,0.95}
\definecolor{red}{rgb}{0.0,0.0,0.0}
\definecolor{green}{rgb}{0.0,0.6,0.0}
\definecolor{blue}{rgb}{0.0,0.0,1.0}
\lstloadlanguages{bash,Java,C,C++,csh,make,sh}%%[Visual]Basic,xml}
\lstset{frame=none,basicstyle=\footnotesize\ttfamily,breaklines,tabsize=2,captionpos=b,prebreak={\hbox{$\rightarrow$}},postbreak={\hbox{$\hookrightarrow$}},showstringspaces=false,backgroundcolor=\color{grey}\bfseries,keywordstyle=\color{blue},commentstyle=\color{green}\textit,stringstyle=\color{red}\ttfamily,abovecaptionskip=2pt,aboveskip=0pt,belowskip=0pt,belowcaptionskip=0pt}
%\lstset{moredelim=[is][\color{red}]{_+_}{_+_}}

\graphicspath{{./pictures/}}
\makeatletter
\usepackage{babel}
\makeatother
\begin{document}
\title{R\'ealisation du syst\`eme pour Red Pitaya sur SD-card \&
acc\`es mat\'eriel en espace utilisateur}

\lhead{}
\chead{Tutorial}
\rhead{Page \thepage/\pageref{LastPage}}
\lfoot{J.-M Friedt}
\cfoot{Universit\'e de Franche-Comt\'e}
%\rfoot{}

\thispagestyle{fancy}

\lstset{numbers=left, numbersep=6pt}
\begin{center}
{\bf \Large
{Buildroot-based Red Pitaya (STEMLab) system development}
}

G. Goavec-M\'erou, J.-M Friedt, \today
\end{center}

Developing an operating system based on the Linux kernel requires
options tailored to the targeted architecture: we will not use, in 
the proposed approach, a pre-build binary package based distribution,
but aim at providing a {\em consistent framework} for generating manually
all tools from sources. The consistent {\em bootloader}, {\em toolchain},
operating system {\em kernel} and filesystem including {\em binary
executables} and {\em libraries} will be provided by Buildroot
\footnote{\url{https://buildroot.org/}}. 

\begin{figure}[h!tb]
\begin{center}
\includegraphics[width=0.8\linewidth]{DSC06098_small.jpg}
\end{center}
\caption{The Red Pitaya board. The USB micro-B connector located closest
to the side of the printed circuit board, opposite to the RJ45 Ethernet
connector, is used for providing the power supply. The USB micro-B connector
located in the middle of the board provides a terminal on a virtual serial
port.}
\label{a13pict}
\end{figure}

\section{Prerequisites}

All commands prefixed by ``\#'' indicate commands to be executed as administrator (``root'').
On some distributions (e.g. Ubuntu), executing as ``root'' requires prefixing the command
with {\tt sudo}, assuming the user is part of the ``sudo'' group (check in {\tt /etc/group} of the 
development host computer if the login is documented after the ``sudo'' group).

All commands prefixed by ``\$'' are to be executed as user. Only execute as few commands as
possible as root whose account should only be used during system administration tasks.

Throughout the document the {\bf host} refers to the computer we are working on -- most often
Intel x86 based architecture, providing a functional GNU/Linux environment through a binary 
distribution (Debian, Ubuntu ...). The host provides a comfortable working environment as
opposed to the embedded {\bf target} whose resources are constrained, usually without graphical
interface, and whose performance must be tailored to the application and not the developer needs.
Our objective will be to {\bf cross-compile} on the host an efficient operating system and
development framework for the target since the latter computer architecture is hardly ever
Intel x86 based.

\colorbox{lightgray}{\begin{minipage}{0.9\textwidth}
\# apt-get install build-essential libncurses5-dev u-boot-tools \\
\# apt-get install git libusb-1.0-0-dev pkg-config 
\end{minipage}} 

%\noindent\fbox{\parbox{\linewidth}{{\bf Attention}~: bien qu'\`a l'encontre
%de la s\'eparation entre l'environnement h\^ote et environnement de d\'eveloppement,
%nous aurons besoin d'une biblioth\`eque de d\'eveloppement de python pour satisfaire
%la compilation de uboot dans buildroot. Installer {\tt libpython2.7-dev}.
%}}

\section{Generating all tools from sources: {\tt buildroot}}
\label{br}

Buildroot \cite{ficheux} is an integrated set of tools aimed at
\begin{enumerate}
\item compiling a custom toolchain tailored to the targeted system (compilation result in 
{\tt output/host/usr/bin/}) and most significantly the data bus width (32, 64 bits) and accelerator
options (FPU, SIMD instructions),
\item compile a tailored Linux kernel,
\item compile associated GNU tools and libraries.
\end{enumerate}

Several options were available when considering which development framework to use with
such a powerful processor as the Zynq fitted on the Red Pitaya: the microprocessor is recent
enough that dedicated functionalities of the Linux kernel must be tailored to the targeted
platform, most significantly accessing the FPGA programmable logic (PL) located next to the 
general purpose dual-processor processing system (PS). Xilinx, manufacturer of the Zynq, has
updated a dedicated development branch of the Linux kernel in parallel to the official branch.
At the moment, only the Xilinx branch, named {\tt linux-xlnx}, allows to configure the FPGA, 
using methods which should eventually be integrated into the official kernel. However, the 
second challenge lies in being able to experiment with the real-time Linux kernel extension -- 
Xenomai -- which is only supported by the official kernels. The dilemma hence lies in selecting
an official kernel without FPGA support, or the divergent Xilinx kernel without Xenomai support, 
with an uncertain long term evolution and availability. Hence in this document two toolchain 
compilations are addressed, one using the Xilinx Linux kernel {\bf with FPGA and Devicetree Overlay
(DTBO)} support, and the other one with Xenomai support relying on a mainline kernel but hence
{\bf not supporting the dynamic devicetree update mechanism of the DTBO}.
% Hoping for a future merge of the
%{\tt linux-xlnx} branch with the official kernel, and considering the size of Xilinx which should 
%guarantee a certain continuity of developments, we have chosen to integrate the Linux kernel in 
%{\tt buildroot} and to adapt the modifications needed to support Xenomai on this kernel -- an 
%evolved version of 4.29.0 at the time of the split from the official kernel. The buildroot extension 
%below will therefore include a specific kernel for the Zynq, supporting the FPGA configuration as 
%well as the Xenomai real-time extension, and the associated development tools, excluding the 
%proprietary Xilinx's proprietary tool for configuring FPGAs named Vivado which is to be manually
%installed separately.

\begin{figwindow}[0, r, \includegraphics{schema},{Elements needed for developing and
running an operating system: the tools needed to generate all these elements must be
consistent.\label{sch}}]
The advantage of buildroot \footnotemark, as well as open-embedded \footnotemark, is
to provide an integrated and consistent environment that hides a number of details from the 
developer (Fig. \ref{sch}). The disadvantage of buildroot, as well as open-embedded, is to 
provide an integrated environment that hides a certain number of details from the developer (!). 
We will take advantage of this working environment to compile all the tools needed to work with the 
Red Pitaya without depending on an external binary distribution, and then to add the Linux 
real-time extension Xenomai \cite{xeno}. This extension requires, in addition to an update of 
the kernel, additional tools in user space that buildroot can manage. Using Xenomai on the Red 
Pitaya will be described in detail in the section \ref{xe}.\\
\indent Buildroot, available at \url{https://github.com/buildroot/buildroot}, does not support 
natively the Red Pitaya platform. We must exploit a mechanism for adding such a functionality 
through the environment variable {\tt BR2\_EXTERNAL} which references the absolute directory 
containing the files missing in the official distribution for the support of a
platform. We will therefore download a second repository in addition to the official distribution 
of {\tt buildroot} in order to add the missing features by pointing to {\tt BR2\_EXTERNAL} to
a tree structure consistent with the file organization expected by by {\tt buildroot}.
\end{figwindow}
 \addtocounter{footnote}{-2} %3=n
 \stepcounter{footnote}\footnotetext{\url{buildroot.uclibc.org/}}
 \stepcounter{footnote}\footnotetext{\url{www.openembedded.org}}

\section{Compiling a Linux kernel and image without Xenomai support
(no real time support but FPGA configuration through DTBO)}

We start from a known official version of {\tt buildroot}, rather than the latest update
available on {\tt github}, in order to use a reliable and reproducible version of each tool: \\
\colorbox{lightgray}{\begin{minipage}{0.9\textwidth}
wget https://buildroot.org/downloads/buildroot-2025.05.1.tar.gz
\end{minipage}} 

The default distribution is complemented with additional features for supporting the Red Pitaya:\\
\colorbox{lightgray}{\begin{minipage}{0.9\textwidth}
git clone https://github.com/trabucayre/redpitaya.git
\end{minipage}} 

\noindent
which creates a {\tt redpitaya} directory holding the necessary files. We must tell 
{\tt buildroot} where these additional files are located: this is achieved with the
environment variable {\tt BR2\_EXTERNAL}. Either this variable can be defined with the
absolute path leading to the {\tt redpitaya} directory, or execute the {\tt sourceme.ggm} script
from within the {\tt redpitaya} directory which was just cloned:\\
\colorbox{lightgray}{\begin{minipage}{0.9\textwidth}
cd redpitaya \\
source sourceme.ggm
\end{minipage}} 

Care must be taken to reload the script in each new terminal, the definition of {\tt BR2\_EXTERNAL} 
being only local to the active terminal and not being propagated to other work sessions.

Enter the {\tt buildroot} directory that was just uncompressed and unarchived with\\
\colorbox{lightgray}{\begin{minipage}{0.9\textwidth}
tar zxvf buildroot-2025.05.1.tar.gz
\end{minipage}} 

{\bf If you plan on renaming this directory, do it now since after the Buildroot compilation
process, some links will be broken if the directory name is changed later.}
From this directory, configure {\tt buildroot} targeted towards the Red Pitaya:\\
\colorbox{lightgray}{\begin{minipage}{0.9\textwidth}
cd buildroot-2025.05.1 \\
make redpitaya\_defconfig 
\end{minipage}} 

Then, {\tt make} generates a first functional image by downloading all the necessary sources
and compile. At the end of this work which can take from several tens of minutes to hours, and
requires 5 to 8~GB of storage space, all Red Pitaya binaries can be found in {\tt output/images}.
% add guillemet fermant a la ligne 21 de redpitaya/configs/redpitaya_defconfig

\includegraphics[width=.7cm]{danger} 
The following lines might {\bf damage any storage medium} if the wrong peripheral is selected.
Alway {\bf check} the name of the peripheral associated with the SD card (\verb~dmesg | tail~) 
before launching the {\tt dd} command.

The image resulting from the compilation sequence is transferred to the SD card with\\
\colorbox{lightgray}{\begin{minipage}{0.9\textwidth}
sudo dd if=output/images/sdcard.img of=/dev/sdc
\end{minipage}} 
in which we have on purpose selected the {\tt /dev/sdc} peripheral since it is hardly ever
the appropriate option: most often, the SD card will be called {\tt /dev/sdb} (second hard disk
compatible with the Linux SCSI driver) or {\tt /dev/mmcblk0} (case of an internal SD card reader).

\noindent\fbox{\parbox{\linewidth}{
\includegraphics[width=.7cm]{danger} %\danger 
{\bf Warning}: the content of the SD-card, or any mass storage peripheral accessed through the
last argument of this command, will be {\bf definitely} lost. Always check twice the name of
the target periperal the Buildroot image is written to.}}

Once the image has been transferred to the SD card, we observe two partitions on the mass
storage medium: the first one formatted as VFAT (compatible with Microsoft Windows) holding the
devicetree description, the Linux kernel, and the bootloader, and the second one holding the 
GNU/Linux filesystem. This second partition size is exactly the one provided as argument in
the Buildroot options {\tt .config} for storing the operating system files, and little space
is available for one own additions -- we will see later that mainly we wish to copy bitstreams
for configuring the FPGA. This lack of space can be corrected by mounting the SD-card and change
the second partition size with
\begin{enumerate}
\item delete the partition using {\tt fdisk}: lauch the program on the host computer as
{\tt root} administrator, display the current configuration (which will be needed below) with
{\tt p}, then {\tt d} to erase partition 2,
\item re-create this partition but this time extending to the whole space available on the
card: {\tt n} to create a new partition, {\tt p} for a primary partition, provide the first
sector address as was just displayed earlier (in the screenshot below, this first sector address
is 32769), and accept the default selection to use the whole available space. Conclude with
{\tt w} to store this new configuration,
\item after leaving {\tt fdisk}, run {\tt resize2fs /dev/sdd2} to resize the filesystem to the size
of the partition.
\end{enumerate}

Following all these operations, the GNU/Linux system uses all of the available space, In the following
example, we started from a configuration that looked like
{\footnotesize
\begin{verbatim}
Device     Boot Start     End Sectors  Size Id Type
/dev/sdd1  *        1   32768   32768   16M  c W95 FAT32 (LBA)
/dev/sdd2       32769 1081344 1048576  512M 83 Linux
\end{verbatim}
}
and after changing the second partition size we end up with
{\footnotesize
\begin{verbatim}
Device     Boot Start     End Sectors  Size Id Type
/dev/sdd1  *        1   32768   32768   16M  c W95 FAT32 (LBA)
/dev/sdd2       32769 7796735 7763967  3.7G 83 Linux
\end{verbatim}
}

Asynchronous (RS232-compatible) link emulation over the USB port located in the middle
of the printed circuit board is functional: launch {\tt minicom} after connecting a USB-micro B
cable, tuned to a 115200 baudrate, 8N1. A prompt stating \verb~redpitaya>~ should appear
upon hitting the ENTER key.

Various additional features can be added such as a graphical interface server (X11,
see appendix \ref{VGA}) but are not detailed in the main part of the text but mentioned
in various appendices.

\subsection{Network configuration}

\begin{figwindow}[0,r,\includegraphics[width=9cm]{sch_ip}, {Architecture of a TCP/IP network}]
{\color{red}
The network configuration between host and target is handled by the interface configuration
command {\tt ifconfig} to assign a logic address (Internet Protocol -- IP) to a physical interface 
(defined through its
MAC address) and create a point to point link between computers of a sub-network and their
gateway. This gateway is in charge of transferring data packets towards computers located outside
the sub-network, or even towards the internet if the targeted address is not within the
subnet. The command to define the routing table is {\tt route}: we shall not need to specify these
routing tables manually in the simple case of the point to point communication of the Red Pitaya
with the host computer since GNU/Linux automatically assigns a given subnet to its interface
and properly routes packets within the subnet.\\ \indent
The default network configuration is to communicate through the Ethernet interface, with the
embedded board IP address given by \verb~ifconfig~. 
}
\end{figwindow}

We observe that
{\footnotesize
\begin{verbatim}
redpitaya> ifconfig                                                             
eth0      Link encap:Ethernet  HWaddr 22:3D:C4:D9:DB:CD                         
          inet addr:192.168.2.200  Bcast:192.168.1.255  Mask:255.255.255.0       
...
\end{verbatim}
}
the default IP address is 192.168.2.200, hence a local sub-network
\footnote{the address ranges for private local networks, which are never
routed towards the internet, are described in RFC1918 available at
\url{https://tools.ietf.org/html/rfc1918}}.

Changing a network interface configuration (on PC or on Red Pitaya) is achieved using the
{\tt ifconfig} command. Arguments include the name of the interface and the associated IP address:

{\verb!ifconfig eth0 192.168.2.1!}

Before configuring an interface, the list of all interfaces, whether already configured or not,
is obtained with {\verb!ifconfig -a!}. This command is also used to identify the hardware
MAC address as found in the {\tt ether} field: this information is the one transmitted to the
gateway when packets are routed towards the internet.

{\tt ifconfig} is slowly becoming obsolete to be replaced with {\tt ip}. The same result as
above is achieved with 
\begin{verbatim}
ip address add 192.168.2.1 dev eth0
ip link set dev eth0 up
\end{verbatim}

{\color{red}The consistent network configuration between two computers is validated
using {\tt ping IP\_dest} requesting an answer ({\tt ping}) from the remote computer
with IP address {\tt IP\_dest}.}

{\color{red}Once a network configuration is completed and manually validated, it can be
automated by writing the configuration in {\tt /etc/network/interfaces}. The rules for
configuring the Red Pitaya IP configuration can be added as follows:
\begin{verbatim}
auto eth0
iface eth0 inet static
  address 192.168.2.200
  netmask 255.255.255.0
  gateway 192.168.2.1
\end{verbatim}

This configuration is activated on the Red Pitaya with {\tt ifup eth0}. If the interface
was already configured, it might be needed to cancel this previous configuration using
{\tt ifdown eth0}. If such a configuration is also used on the PC, make sure the Network Manager
is not conflicting with {\tt /etc/network/interfaces}: Network Manager can be told not to mess
with interfaces defined in {\tt /etc/network/interfaces} by stating {\tt managed=false} 
in \\ {\tt /etc/NetworkManager/NetworkManager.conf} of the PC Linux distribution.

%Sur Redpitaya, u-boot se charge de fournir la configuration par d\'efaut de
%l'interface eth0 lors du d\'emarrage de la carte. La configuration de
%u-boot se trouve dans {\tt redpitaya/board/redpitaya/uEnv.txt}~: nous y
%trouvons l'adresse IP {\tt ipaddr} ainsi que les interlocuteurs que sont
%le serveur d'adresses IP s'il y a lieu et la passerelle par d\'efaut
%{\tt gatewayip}.
%
%Modifier {\tt uEnv.txt} afin de configurer par d\'efaut l'adresse IP de
%la Redpitaya \`a 192.168.2.200 et la passerelle par d\'efaut \`a
%192.168.2.1
%
%\begin{verbatim}
%ipaddr=192.168.2.200
%serverip=192.168.2.1
%gatewayip=192.168.2.1
%\end{verbatim}

\label{lighttpd}
We shall later need a web server: {\tt lighttpd} is a light and stable option.
Adding a new application with Buildroot is achieved by running
{\tt make menuconfig} and search for the location where the web server can be
activated, possibly by searching (\verb~/~ followed by {\tt lightt}) a keyword, 
or by activating directly {\tt Target packages} $\rightarrow$
{\tt Networking applications} $\rightarrow$ {\tt lighttpd}. 

Life will be easier by installing a text editor on the target system. The most
efficient, {\tt vim}, requires activating the {\tt wchar} and {\tt curses} of
the toolchain with %the sequence {\tt Toolchain} $\rightarrow$ 
%{\tt Enable WCHAR support} and 
{\tt Target packages} $\rightarrow$
{\tt Libraries} $\rightarrow$ {\tt Text and terminal handling} 
$\rightarrow$ {\tt ncurses} and {\tt enable wide char support} of {\tt ncurses}. Once these libraries have been activated,
{\tt vim} support in {\tt busybox} is activated by {\tt Target packages} $\rightarrow$ 
{\tt Show packages that are also provided by busybox} and {\tt Target packages} 
$\rightarrow$ {\tt Text editors and viewers} $\rightarrow$ {\tt vim}. Many other alternative
text editors are also available in the same menu. Once selected, quit {\tt make menuconfig},
run {\tt make} and transfer the resulting image to the SD card upon completion of the compilation.

The addition of the new features is validated by launching the Red Pitaya and
making sure the new software is indeed available:
\begin{verbatim}
redpitaya> which lighttpd                                                       
/usr/sbin/lighttpd
\end{verbatim}

\subsubsection{Domain name correspondence with IP addresses}

{\color{red}
The domain name server, in charge of translating the domain name to an IP address -- DNS --
is defined in {\tt /etc/resolv.conf}. This information is needed only for reaching
domains on the internet, for example {\tt free.fr} or {\tt google.com}, but is useless
as long as we are restricted to local connections in which the targeted system is
identified with its IP address. The DNS information is already filled on the PCs, but the
Red Pitaya is unaware of the DNS IP address to relate a domain name, e.g. {\tt sequanux.org},
and its IP address (188.165.36.56). Filling the configuration file \\{\tt /etc/resolv.conf} 
with information such as

{\tt nameserver 194.57.91.200 \\
nameserver 194.57.91.201}\\
with 194.47.91.20\{0,1\} the domain name servers of Universit\'e de 
Franche Comt\'e will solve the issue. On a different site, the DNS address will be found
by looking into the {\tt /etc/resolv.conf} of a local PC. A generic solution is
to selected {\tt 9.9.9.9} as DNS\footnote{\url{https://www.quad9.net/}}.
server address
}

When connecting using the Secure Shell {\tt ssh} from a PC to the Red Pitaya, the only administrator
account {\tt root} password must be provided.

\fbox{\parbox{\linewidth}{The default {\tt root} password is {\bf root}.}}

\subsubsection{NFS ({\em Network File System})}\label{nfsexport}

In the case of room 235B, a directory has been allowed on each PC to be shared with the
Red Pitaya: this directory is {\tt /home/etudiant/nfs}. This directory represents the more generic
reference to {\tt /home/utilisateur/target} in the general description that follows.

Sharing a directory by NFS allows for making the user believe that part of the (host) PC hard
disk drive can be accessed as a directory on the (target) embedded Red Pitaya board. Sharing
virtually files through such a network connection provides a flexible working environment since
after cross-compiling a binary file on the (fast) PC, the resulting executable is immediately
available in the shared directory for executing on the embedded board.

Achieving this result requires, on the PC, to create a {\tt /home/utilisateur/target} 
directory which will be the shared directory. Then in the {\tt /etc/exports} configuration
file of the host (running the NFS server), we allow sharing this directory by adding
{\tt /home/utilisateur/target *(rw,sync)}

Restart the NFS server on the PC to comply with the new rule:\\
\colorbox{lightgray}{\begin{minipage}{0.9\textwidth}
\#/etc/init.d/nfs-kernel-server restart
\end{minipage}}\\

This directory will allow sharing files between the host and the target, for examples
executables cross-compiled on the PC and we wish to execute on the Red Pitaya, or in the
other direction data that were collected on the Red Pitaya to be processed on the host
using a graphical interface such as {\tt gnuplot} or GNU/Octave for displaying the charts.
For accessing the directory on the Red Pitaya, the remote directory is mounted using
{\tt mount -o nolock IP\_PC:/home/utilisateur/target /mnt}: the content of {\tt 
/home/utilisateur/target} on the host is now accessible in the {\tt /mnt} directory of the 
target. In room 235B, this result is achieved by executing on the Red Pitaya\\
{\tt mount -o nolock 192.168.2.1:/home/etudiant/nfs /mnt} \\
assuming the Red Pitaya IP address is 192.168.2.x (on the same sub-network than the PC). If the protocol
selection fails, the {\tt -t nfs4} option might be added.

\textbf{Connecting the Red Pitaya to the internet by configuring the PC as gateway (optional)}

We only briefly mention the ability to use the PC as a gateway that translates the
Red Pitaya IP requests and broadcast towards the targeted server on the internet: this
capability will not be used at first and can be skipped if not needed. The tool for
achieving such address translation is called {\tt iptables}: a configuration example
for transferring all packets between {\tt eth0} and {\tt eth1} interfaces of the host PC is
provided below. {\color{red}This configuration only makes sense on the PC where we assume
that {\tt eth0} is connected to the subnetwork linked to the gateway allowing to reach the
internet, and {\tt eth1} is connected to the subnetwork holding the Red Pitaya. These commands
all require administration ({\tt root}) credentials.}

\begin{verbatim}
monpath=/sbin
echo "1" > /proc/sys/net/ipv4/ip_forward
echo "1" > /proc/sys/net/ipv4/ip_dynaddr
${monpath}/iptables -P INPUT ACCEPT
${monpath}/iptables -F INPUT 
${monpath}/iptables -P OUTPUT ACCEPT
${monpath}/iptables -F OUTPUT 
${monpath}/iptables -P FORWARD ACCEPT
${monpath}/iptables -F FORWARD
${monpath}/iptables -t nat -F

${monpath}/iptables -A FORWARD -i eth0 -o eth1 -m state --state ESTABLISHED,RELATED -j ACCEPT
${monpath}/iptables -A FORWARD -i eth1 -o eth0 -j ACCEPT
${monpath}/iptables -A FORWARD -j LOG

${monpath}/iptables -t nat -A POSTROUTING -o eth0 -j MASQUERADE
\end{verbatim}

\subsection{Executing in an emulator}

Lacking hardware as is often the case for software developers of embedded systems, the
image resulting from {\tt buildroot} compilation can be executed in an emulator, here
{\tt qemu} emulating an ARM processor. We will use {\color{red}at first, as long as the {\tt eth0}
network interface is not needed}, the Debian GNU/Linux package version provided as {\tt 
qemu-system-arm}. Once this packet has been installed, go to the Buildroot {\tt output/images} 
directory and run
\begin{verbatim}
qemu-system-arm -append "console=ttyPS0,115200 root=/dev/mmcblk0 rw earlyprintk" -M xilinx-zynq-a9 \
-m 1024  -serial mon:stdio -dtb zynq-red_pitaya.dtb -nographic -kernel zImage  -sd rootfs.ext4
\end{verbatim}
to boot Linux, load the ({\tt rootfs}) filesystem and complete loading the system until a prompt
appears. All system commands are functional, but accessing hardware requires of course that
such peripherals are specifically emulated, a functionality which {\tt qemu} might not provide.
For example, GPIO related registers seem to be properly emulated
\begin{verbatim}
redpitaya> echo "908" > /sys/class/gpio/export 
redpitaya> echo "out"  > /sys/class/gpio/gpio908/direction 
redpitaya> echo "1"  > /sys/class/gpio/gpio908/value
\end{verbatim}
even though the consequences of writing to a peripheral controling register are not obvious.

The drawbacks of this approach are that 1/ the Ethernet interface of the Red Pitaya is not
emulated and 2/ only the guest operating system on the Red Pitaya emulated with 
{\tt qemu} can access the host operating system on the PC running {\tt qemu}, but connections
from the host to the guest are are not supported as long as dedicated network interface has
not been created. In order to compensate for these limitations, the specific branch of
{\tt qemu} provided by Xilinx as described at 
\url{https://xilinx-wiki.atlassian.net/wiki/spaces/A/pages/821395464/QEMU+User+Documentation}
% https://xilinx-wiki.atlassian.net/wiki/spaces/A/pages/18842060/QEMU} 
is compiled, and most significantly with the support of the Red Pitaya Zynq processor as
detailed at \url{https://xilinx-wiki.atlassian.net/wiki/spaces/A/pages/189530183/Zynq-7000}.
% \url{https://xilinx-wiki.atlassian.net/wiki/spaces/A/pages/18842054/QEMU+-+Zynq-7000}. 
Notice that compiling all of {\tt qemu} is not needed but that after adding the devicetree
compiler with \verb~git submodule update --init dtc~ we will only compile the wanted target,
namely \\
% 260327 \verb~$./configure --target-list="aarch64-softmmu" --enable-fdt --disable-kvm --disable-xen~
\verb~$./configure --target-list="arm-softmmu" --disable-werror --enable-fdt --disable-kvm --disable-xen~

Thanks to the compilation of the binary resulting from the archive provided by Xilinx at
\url{https://github.com/Xilinx/qemu} as described below, we end up with 
{\tt qemu-system-arm} with the support of the {\tt arm-generic-fdt-7series}
machine whose Ethernet interface is detected as {\tt eth0}. The command line to launch 
{\tt qemu} with such a functionality is
%\begin{verbatim}
%sudo .../aarch64-softmmu/qemu-system-aarch64 -append "console=ttyPS0,115200 root=/dev/mmcblk0 \
%rw earlyprintk" -M arm-generic-fdt-7series -machine linux=on -smp 2 -m 1024  -serial mon:stdio \
%-dtb zynq-red_pitaya.dtb -nographic -kernel zImage  -sd rootfs.ext4  -net nic -net nic -net nic \
%-net nic -net tap,downscript=no
%\end{verbatim}
\begin{verbatim}
sudo .../qemu-system-arm -append "console=ttyPS0,115200 root=/dev/mmcblk0 -boot mode=5 \
rw earlyprintk" -M arm-generic-fdt-7series -machine linux=on -smp 2 -m 1024  -serial mon:stdio \
-dtb zynq-red_pitaya.dtb -nographic -kernel zImage -initrd rootfs.cpio.gz \
-net nic -net nic -net nic -net nic -net tap,downscript=no
\end{verbatim}

\noindent
with the whole command launched from Buildroot's {\tt output/images} directory to avoid filling
the full path of all the files (devicetree {\tt zynq-red\_pitaya.dtb}, root filesystem 
{\tt rootfs.ext4}, Linux kernel {\tt zImage}). Notice that this command must be launched as
{\bf administrator (root) on the PC} in order to be allowed to create the {\tt tap0} network 
interface on the host.

{\color{red}Once the virtual Red Pitaya launched in {\tt qemu-system-aarch64}, we observe with
({\tt ifconfig -a}) that the {\tt eth0} network interface on the target is available. On the PC, 
a new interface was created when launching {\tt qemu} named {\tt tap0}: on the PC, {\tt sudo 
ifconfig  -a | grep tap} displays the properties of this new interface. {\bf Beware}: if
an older execution of qemu was improperly killed or is running as a background task, a new {\tt tap} 
interface such as {\tt tap1} might be created, Make sure to properly identify which interface
is connected to which running copy of {\tt qemu}.

The {\tt tap} interface having been identified, we must still configure an address on the PC in
the same subnet -- but of course with a different least significant byte -- than the Red Pitaya 
address. Once this configuration is completed on the host and the target, a successful {\tt ping} 
validates the proper operation of the link.
}

Figure \ref{rpqemu} demonstrates the configuration of the {\tt eth0} interface when emulating the
Red Pitaya (address 192.168.0.10) in {\tt qemu}, consistent with the {\tt tap0} interface created
and configured on the host (same sub-net 192.168.0.x with the PC arbitrarily assigned the 
192.168.0.5 address) and the ability of a web client running on the host to access the
{\tt lighttpd} server executed on the guest.

\begin{figure}[h!tb]
\includegraphics[width=\linewidth]{redpitaya_qemu.png}
\caption{Top-left the guest: Xilinx {\tt qemu} emulating the Red Pitaya. Bottom-left the web 
client executed on the host and connected to the web server running on the guest. Bottom-right the
host network configuration created when launching {\tt qemu} with the {\tt -net tap} option.}
\label{rpqemu}
\end{figure}

A last advice: if we wish to avoid the hassle of secure shell connexion as provided by
{\tt ssh} between the host and the guest, a telnet server {\tt telnetd} is available with
{\tt buildroot} for the Red Pitaya as a feature provided by busybox accessed with
\verb~make busybox-menuconfig~$\rightarrow$\verb~networking~$\rightarrow$\verb~telnetd~. 
With this server, the Red Pitaya becomes accessible with {\tt telnet 192.168.0.10} (according
to the configuration visible on Fig. \ref{rpqemu}) without performance degradation due to
sharing public cryptographic keys, but with information circulating un-encrypted over the
network.

Files can be transferred from host to target using {\tt scp}. Furthermore, NFS is also functional
in this framework, as long as a functional TCP/IP link exists between host and guest and that
the host directory to be shared has been defined by including the guest IP address in
{\tt /etc/exports} (see previously section \ref{nfsexport}).

Leaving {\tt qemu} cleanly while resetting all variable and memory allocations is achieved
similar {\tt minicom}, with ``CTRL-a'' followed by ``x''.
 
\subsection{Web server}\label{boa}

We have seen (section \ref{lighttpd}) how to complement the Redpitaya Buildroot configuration in order
to add support for a light web server, implementing CGI script execution, named
{\tt lighttpd}. A default configuration is found in {\tt /etc/lighttpd/lighttpd.conf} 
using the {\tt /var/www/} directory for storing the files accessible through a web
connection. The server is launched as a daemon using 
{\tt lighttpd -f /etc/lighttpd/lighttpd.conf}.
If an error occurs upon launching {\tt lighttpd}, check the error log located in
{\tt /var/log/lighttpd-error.log} that the issue is not related to
``pcre support is missing for: dir-listing.exclude''. In case this error occurs, we 
might deactivate the faulty module by commenting out
\begin{verbatim}
# include "conf.d/dirlisting.conf"
\end{verbatim}
in {\tt /etc/lighttpd/lighttpd.conf}.

A basic static web page allows for testing proper operation of the web server and connectivity
to the web browser on the host computer: on the Red Pitaya, edit a file named {\tt index.html}
in the {\tt /var/www} directory including the sentence {\tt Hello World} and check that it
is displayed when typing the Red Pitaya IP address in the host web browser.

Further interaction with hardware will require the ability to dynamically analyze the URL arguments
and act accordingly, hence launching a script on the Red Pitaya when requested by the host web browser.
Indeed, CGI ({\em Common Gateway Interface}) allows for dynamically generating web pages
by executing a script generating HTML sentences that can be read by the client. The protocol
specification, \url{http://www.ietf.org/rfc/rfc3875.txt}, indicates the means for sharing information
between client and server.

CGI script support is added by storing in {\tt /etc/lighttpd/conf.d/cgi.conf} a configuration
file with the following content
\begin{verbatim}
server.modules += ( "mod_cgi" )
cgi.assign                 = ( ".pl"  => "/usr/bin/perl",
                               ".cgi" => "/bin/sh",
                               ".py"  => "/usr/bin/python",
                               ".php" => "/usr/bin/php-cgi" )

index-file.names           += ( "index.pl",   "default.pl",
                               "index.rb",   "default.rb",
                               "index.erb",  "default.erb",
                               "index.py",   "default.py",
                               "index.php",  "default.php" )
\end{verbatim}

and adding {\bf at the end} of {\tt /etc/lighttpd/lighttpd.conf} the line 
{\tt include "conf.d/cgi.conf"} for loading this CGI configuration file.

Once the update is completed, the web server is restarted by using
{\tt /etc/init.d/S50lighttpd restart}, possibly after killing the previous process
that was launched manually by finding its identifier with the {\tt ps} command. 


%\begin{enumerate}
%\item {\tt mkdir /tmp/boa}
%\item {\tt chown nobody.nogroup /tmp/boa}
%\item {\tt chmod 755 /tmp/boa} (v\'erifier que l'utilisateur nobody.nogroup peut acc\'eder \`a ce r\'epertoire)
%\item placer dans {\tt /tmp/boa/} un fichier de configuration {\tt boa.conf} contenant par exemple
%\begin{verbatim}
%Port 80
%User nobody
%Group nogroup
%ErrorLog /tmp/boa/error_log
%AccessLog /tmp/boa/access_log
%DocumentRoot /tmp/boa/
%UserDir public_html
%DirectoryIndex index.html
%DirectoryMaker /tmp/boa/boa_indexer
%# DirectoryCache /var/spool/boa/dircache
%KeepAliveMax 1000
%KeepAliveTimeout 10
%MimeTypes /etc/mime.types
%DefaultType text/plain
%# Uncomment the next line if you want .cgi files to execute from anywhere
%AddType application/x-httpd-cgi cgi
%ScriptAlias /cgi-bin/ /tmp/boa/cgi/
%\end{verbatim}

\fbox{\parbox{\linewidth}{
\includegraphics[width=.7cm]{danger} %\danger 
Never run services with administration privileges ({\tt root}). Even if such 
a configuration makes accessing hardware resources much easier, it break Unix security
rules based on the clear separation between services (examples of poor web server
configuration exploits: {\em Exploiting Network Surveillance Cameras Like a Hollywood Hacker}, 
Black Hat 2013, at \url{https://www.youtube.com/watch?v=B8DjTcANBx0}).  Here, the web server
is run as user and group {\tt www-data} which have no administration rights so that a security
failure in the web server would not lead to a security breach in the whole system.}}
%\item lancer {\tt boa} par {\tt boa -f /tmp/boa/boa.conf}
%\item connecter un client web vers l'adresse IP de la carte Red Pitaya~: le fichier
%{\tt index.html} que nous aurons plac\'e dans {\tt /tmp/boa} y sera affiche\'e si les
%permissions appropri\'ees lui sont accord\'ees (prendre soin \`a affecter le bon
%propri\'etaire \`a ce fichier s'il est cr\'e\'e par {\tt root}).
%\end{enumerate}

For executing CGI scripts, we might create a sub-directory {\tt /var/www/cgi} (in agreement
with the configuration file) and store there the shell script (since {\tt .cgi} extensions
are associated with bash scripts):
\begin{verbatim}
#!/bin/sh
echo -e "Content-type: text/html\r\n\r\n"
echo "Hello CGI" $(date)
\end{verbatim}

which is launched from the web client by providing the url {\tt http://IP/cgi/script.cgi}. 
We will of course remember to make the script executable ({\tt chmod 755 script.cgi} in the {\tt cgi} directory). This method hence
allows for dynamically generating web pages, for example for displaying sensor reading measurements, or in
this case the current date.

The script running on the server can, for example, read a sensor or GPIO state and update the web
page accordingly. More interesting, the HTML request can include variable names and their values
to be used by the script. The two methods meeting this objective are POST and GET. The latter will
be used here since the data are shared between client and server through the url, easier to implement
client and server \footnote{notice that the apparent protection provided by the POST post method
against a malicious usage is trivially bypassed with tools such as {\tt curl}}. Hence, providing arguments
to the CGI script is achieved with URLs stating
\begin{verbatim}
http://192.168.2.200/cgi/script.cgi?variable1=val1&variable2=val2&var3=val
\end{verbatim}

The only subtely in this usage is
\begin{enumerate}
\item understand the execution context of the CGI script, whatever language it is written in
(the first line, as in all scripts, informs on the location of the interpreter),
\item fetch the list of variables provided through the url. The variables including these
parameters are described at page 8 of \url{http://www.ietf.org/rfc/rfc3875.txt}.
\end{enumerate}

Before we can interact with hardware from a CGI script, we must learn how to interact with
GPIOs from the interactive shell on the Red Pitaya.

\subsection{Accessing hardware from the shell}\label{acces_led}

If we wish to use the LED control driver: the {\tt /sys/class/leds} virtual directory
provides access to the associated kernel driver for controlling the LEDs. The following sequence
{\footnotesize
\begin{verbatim}
redpitaya> cd /sys/class/leds/
redpitaya> echo "1" > led8/brightness   # orange
redpitaya> echo "0" > led8/brightness 
redpitaya> echo "0" > led9/brightness   # rouge
redpitaya> echo "1" > led9/brightness 
\end{verbatim}
}
allows for checking proper hardware resource control by the Linux kernel.

The {\tt leds} interface is more limited than the general {\tt gpio} interface but 
easier to use. Various interfaces are provided to the user to access hardware resources.
The current trend is to expose these interfaces through the {\tt /sys} directory. 
Two methods are provided for accessing GPIOs connected to the processor (MIO in the Zynq
naming convention, as opposed to EMIO connected to the PL) thanks to the Linux drivers: 
{\tt /sys/class/leds} and {\tt /sys/class/gpio}.

A complementary approach is to use {\tt /sys/class/gpio}
\footnote{\url{https://www.thegoodpenguin.co.uk/blog/stop-using-sys-class-gpio-its-deprecated/} claims
that {\tt /sys/class/gpio} is becoming obsolete, but remains available in the 6.6 kernel so far ...} which provides lower access
to the GPIO since LEDs are a sub-set of GPIOs: the LED driver relies on the GPIO driver.

If drivers are compiled as modules, we observe that they are loaded with
{\footnotesize
\begin{verbatim}
redpitaya> lsmod
Module                  Size  Used by    Not tainted
gpio_zynq               6509  2 
leds_gpio               2890  0 
led_class               3299  1 leds_gpio
\end{verbatim}
}

For accessing GPIOs, a driver implements {\tt /sys/class/gpio}
\footnote{make sure to activate such functionalities in the kernel configuration with
{\tt Device Drivers}$\rightarrow${\tt GPIO Support}$\rightarrow${\tt /sys/class/gpio/... 
(sysfs interface)}}. One feature of {\tt /sys} is to communicate through human readable 
sentences (as opposed to binary values shared with peripherals accessed through {\tt /dev}).

\begin{framed}
\noindent{\bf Kernel versions $\geq6$ (v.s 4., 5.x)}

Practically, the {\tt gpiolib} drivers assumes that a digital port exposes at most
32 input-output pins. Hence, the index of pin {\tt n} on port {\tt P} is 
{\tt P}$\times 32+${\tt n}, with port A indexed as 0, B indexed as 1 and so on. In the case
of the Red Pitaya, pins referred to as MIO are indexed starting with the value 906 in kernels
prior to 6, and with values starting with 512 for kernels 6.x. Hence,
MIO0 is indexed as 906 on pre-6 kernel and 512 on the 6.x kernel, and MIO7 is indexed as 913 on
pre-6 kernel and as 519 on 6.x kernels: the two LEDs -- orange and red -- are connected
to these two pins. We demonstrate control of these peripherals with the following sequence.

In a first step, we request control of the resource with 
{\footnotesize
\begin{verbatim}
redpitaya> echo "906" > /sys/class/gpio/export # for kernel < 6, 512+ otherwise
-sh: write error: Device or resource busy
\end{verbatim}
}
which fails since the LED driver already controls these resources, emphasizing the role of the 
kernel in maintaining consistency when accessing a given resource by multiple users.
\end{framed}


Since kernel 6 (check with {\tt uname -a} which kernel version you are running), the GPIO indices
have been renamed as can be checked with the {\tt debugfs} available when compiling the Linux
kernel with {\tt DEBUG\_FS=y} and {\tt mount -t debugfs none /sys/kernel/debug/} so that
{\tt cat /sys/kernel/debug/gpio} displays

\begin{verbatim}
gpiochip0: GPIOs 512-629, parent: platform/e000a000.gpio, zynq_gpio: 
 gpio-512 (                    |led8                ) out lo 
 gpio-519 (                    |led9                ) out lo  
\end{verbatim}

As in the previous discussion, but replacing 906 or 913 with 512 and 519 respectively, we can
access the GPIO with

\begin{verbatim}
echo 512 > /sys/class/gpio/export
echo out > /sys/class/gpio/gpio512/direction   
echo 1 > /sys/class/gpio/gpio512/value 
\end{verbatim}

\noindent only once the LED driver has freed the resources.

\begin{figure}[h!tb]
\includegraphics[width=0.75\linewidth]{2026-03-14-130029_2704x1050_scrot.png}
\includegraphics[width=0.24\linewidth]{led_sch}
 % led_doc.png} 
\caption{LEDs available on the Red Pitaya: notice that only LED8 and LED9 are accessible
from the PS (i.e. without configuring the FPGA of the PL). Figures taken from 
{\tt https://redpitaya.readthedocs.io/en/latest/developerGuide/fpga/advanced/signal\_mapping.html} and
% {\tt media.readthedocs.org/pdf/rpdocs/latest/rpdocs.pdf} and
{\tt downloads.redpitaya.com/doc/Red\_Pitaya\_Schematics\_STEM\_125-10\_V1.0.pdf}
respectively}.
\label{led}
\end{figure}

We have taken care to configure the drivers handling hardware resources as modules, so that
these resources can be freed by removing the drivers from memory:

{\footnotesize
\begin{verbatim}
redpitaya> rmmod leds_gpio
redpitaya> rmmod led_class
redpitaya> echo "906" > /sys/class/gpio/export  # for pre-6 kernel, 512 otherwise
\end{verbatim}
}

\noindent for pre-6 kernels, or for 6.x kernels:

{\footnotesize
\begin{verbatim}
redpitaya> rmmod leds_gpio
redpitaya> rmmod led_class
redpitaya> echo "512" > /sys/class/gpio/export 
\end{verbatim}
}

Now, the previous request succeeds -- no error message is given -- and leads to the creation
of a new directory holding the pseudo-files for accessing the hardware resources
through requests to the kernel: {\tt gpio906} for pre-6 kernels and  {\tt gpio512} for 6.x kernels,
\footnote{
% \url{http://forum.redpitaya.com/viewtopic.php?f=14&t=115}
\url{https://forum.redpitaya.com/viewtopic.php?t=1158}
explains that LED8 = MIO[0]=gpio906 and LED9 = MIO[7]=gpio913}. These pseudo-files handle the
communication direction ({\tt out} or {\tt in}) and, for pins configured as outputs, their state
({\tt 1} or {\tt 0}).

{\footnotesize
\begin{verbatim}
redpitaya> echo "out" > /sys/class/gpio/gpio906/direction 
redpitaya> echo "1" > /sys/class/gpio/gpio906/value 
redpitaya> echo "0" > /sys/class/gpio/gpio906/value 
\end{verbatim}
}

\noindent for pre-6 kernels and

{\footnotesize
\begin{verbatim}
redpitaya> echo "out" > /sys/class/gpio/gpio512/direction 
redpitaya> echo "1" > /sys/class/gpio/gpio512/value 
redpitaya> echo "0" > /sys/class/gpio/gpio512/value 
\end{verbatim}
}

\noindent for 6.x kernels, validates that accessing hardware resources from the shell is functional.

\subsection{Back to lighttpd and CGI scripts}

Now that we understand access to the GPIO and LED hardware from the shell, this knowledge
is transposed to the CGI script. Be aware though that the main challenge will be ownership
of the resources and being allowed to access hardware as non-administrator user.

A basic example of a functional CGI script is provided below, which requires to understand
the last line to refer to the section concerning hardware access at \ref{acces_led}, as
% echo 12 > /sys/class/gpio/export 
% echo out > /sys/class/gpio/gpio12_pg9/direction
\begin{verbatim}
[root@buildroot cgi]# less cgi/script.cgi 
#!/bin/sh
echo -e "Content-type: text/html\r\n\r\n"
echo "Hello CGI<br>"
echo $QUERY_STRING "<br>"
led=`echo $QUERY_STRING | cut -d= -f2`
echo $led "<br>"
echo $led > /sys/class/leds/led8/brightness
\end{verbatim}
located in {\tt /var/www/cgi}.
% (tel que d\'efini dans la variable {\tt ScriptAlias} du fichier de configuration de {\tt boa}). 
The string provided as argument in the URL is accessed through the variable \verb~$QUERY_STRING~ 
which is here split in its individual items by using the {\tt cut} command.

We must however make sure that the owner of the web server, defined as the
user {\tt www-data} for {\tt lighttpd} to isolate permissions and limit the
risks of malicious access in case security of the server is breached, has access
to the requested hardware resources. That is usually not the case and we must
manually update permissions:
% \begin{verbatim}
% [nobody@buildroot ~]$ echo "1" > /sys/class/gpio/gpio12_pg9/value 
% -sh: /sys/class/gpio/gpio12_pg9/value: Permission denied
% [nobody@buildroot ~]$ logout
% [root@buildroot cgi]# chmod 777 /sys/class/gpio/gpio12_pg9/value 
% [root@buildroot cgi]# su - nobody
% [nobody@buildroot ~]$ echo "0" > /sys/class/gpio/gpio12_pg9/value 
% [nobody@buildroot ~]$ echo "1" > /sys/class/gpio/gpio12_pg9/value
\begin{verbatim}
# ls -l /sys/class/leds/led*/*
lrwxrwxrwx    1 root     root             0 Jan  1 00:13 brightness
...
# chown www-data:www-data /sys/class/leds/led*/brightness
# ls -l /sys/class/leds/led*/
\end{verbatim} 
for {\tt lighttpd}.

\noindent{\bf Demonstrate how you can define the state of {\em both} LEDs 8 and 9 by providing
a GET url to {\tt lighttpd} with both arguments}.

% echo '<meta http-equiv="Refresh" content="3; url=http://192.168.1.172:/timedomain.html"/>'

\subsection{Concept of {\tt devicetree}}

A SOC ({\em System On Chip}) provides many peripherals around a same processing core.
Describing these peripherals can become complex and require inheriting properties: on
ARM architectures, the most recent approach lies in using the {\tt devicetree} as
was already done on SPARC or PowerPC architectures. This file, with the extension
{\tt dts} to be readable by a programmer, is compiled as a {\tt dtb} binary file readable
by kernel modules and drivers.

Accessing LEDs through the {\tt /sys/class/leds} interface configured by the
{\tt devicetree} is achieved by declaring which pin is associated with the functionality,
leading to the creation of a directory holding the pseudo-files for handling pin states.

We observe the configuration of MIO0 and MIO7 as interfaces led8 and led9 respectively
in the devicetree section found in {\tt redpitaya/board/redpitaya/patches/linux/xilinx-v2024.1/0002-redpitaya12-add\_devicetree.patch}
% redpitaya/board/redpitaya/zynq-red\_pitaya.dts} 
of \\ {\tt BR2\_EXTERNAL} defining the Red Pitaya peripherals:

{\footnotesize
\begin{verbatim}
gpio-leds {
 compatible = "gpio-leds";
 led-8-yellow {
  label = "led8";
  gpios = <&gpio0 0 0>;
  default-state = "off";
  linux,default-trigger = "mmc0";
 };
 led-9-red {
  label = "led9";
  gpios = <&gpio0 7 0>;
  default-state = "off";
  linux,default-trigger = "heartbeat";
 };
};
\end{verbatim}
}

The {\em devicetree} will be our most trusted interface when defining specific platform configurations
to the Linux kernel, as will be described later in detail, especially when configuring
the programmable logic (PL) of the chip.

\section{Cross-compiling using the Buildroot toolchain}

C and C++ programs are cross-compiled using the Buildroot toolchain using the tools available in the
{\tt output/host/usr/bin} subdirectory. Make sure to add this path to \verb~$PATH~ in order to use
{\tt arm-linux-gcc} for 32-bit ARM (as is the case for Red Pitaya), or {\tt aarch64-linux-gcc} for 64-bit
ARM (e.g. for Raspberry Pi 4 or 5).

\subsection{TCP/IP and UDP/IP servers using a C program}

While Python3 will provide compact and easy means for implementing a server on the Red Pitaya, the
framework is complex to install and not readily available in our Buildroot image. We shall hence focus
on a C implementation of a TCP and UDP server over IP and check how they interact from the host
using {\tt telnet} for accessing the TCP server or {\tt netcat} for accessing the UDP server.

\begin{lstlisting}[language=C,basicstyle=\footnotesize\ttfamily]
#include <sys/socket.h>
#include <resolv.h>
#include <unistd.h>
#include <strings.h>
#include <arpa/inet.h>

#define MY_PORT         9999
#define MAXBUF          1024

int main()
{int sockfd;
 struct sockaddr_in self;
 char buffer[MAXBUF];
 // socket type (AF = IPv4, STREAM=TCP)
 sockfd = socket(AF_INET, SOCK_STREAM, 0); 
 bzero(&self, sizeof(self));
 self.sin_family = AF_INET;
 self.sin_port = htons(MY_PORT);
 self.sin_addr.s_addr = INADDR_ANY;
 bind(sockfd, (struct sockaddr*)&self, sizeof(self));
 listen(sockfd, 20);                        // wait for incoming connection
 while (1)
 {struct sockaddr_in client_addr;
  int mysize,clientfd;
  unsigned int addrlen=sizeof(client_addr);
  clientfd = accept(sockfd, (struct sockaddr*)&client_addr, &addrlen);
  printf("%s:%d connected\n", inet_ntoa(client_addr.sin_addr), ntohs(client_addr.sin_port));
  mysize=recv(clientfd, buffer, MAXBUF, 0);
  send(clientfd, buffer, mysize, 0);
  close(clientfd);
 }
 close(sockfd);return(0);  // Clean up (should never get here)
}
\end{lstlisting}
implements the TCP server sequence
\begin{enumerate}
\item socket (protocol) 
\item bind (port) 
\item listen (blocking wait)
\item accept 
\item send/recv 
\item close
\end{enumerate}
Be aware of the endianness of both speakers when sharing multi-byte data structures: {\tt hton} (host to network)
and {\tt ntoh} (network to host) followed by {\tt s} for short (two bytes) or {\tt l} (long for 4 bytes) make sure
the endianness between both speakers is consistent: the internet is by convention {\bf big-endian}, while most processor
architectures (x86, most ARM, RISC-V) are little-endian, hence requiring the conversion.

For a UDP server:
\begin{lstlisting}[language=C,basicstyle=\footnotesize\ttfamily]
#include <unistd.h>
#include <strings.h>
#include <sys/socket.h>
#include <resolv.h>
#include <arpa/inet.h>

#define BUFSIZE         1024

void alltoupper(char* s)
{while ( *s != 0 ) *s++ = toupper(*s);}

int main()
{ char buffer[BUFSIZE];
  struct sockaddr_in addr;
  int sd, addr_size, bytes_read;

  sd = socket(PF_INET, SOCK_DGRAM, 0);
  addr.sin_family = AF_INET;
  addr.sin_port = htons(9999);
  addr.sin_addr.s_addr = INADDR_ANY;
  bind(sd, (struct sockaddr*)&addr, sizeof(addr));
  do {bzero(buffer, BUFSIZE);addr_size = BUFSIZE;
      bytes_read=recvfrom(sd,buffer,BUFSIZE,0, \
         (struct sockaddr*)&addr,&addr_size);
      printf("Connect: %s:%d %s\n",inet_ntoa(addr.sin_addr),\
         ntohs(addr.sin_port), buffer);
      alltoupper(buffer);
      sendto(sd,buffer,bytes_read,0,(struct sockaddr*)&addr, \
         addr_size);
  } while ( bytes_read > 0 );
  close(sd);return 0;
}
\end{lstlisting}

Notice how the {\tt socket()} argument was switched from {\tt SOCK\_STREAM} to {\tt SOCK\_DGRAM}. For
testing, {\tt netcat} is used with the {\tt -u} argument for UDP followed by the IP and port: 
{\tt echo "toto" | nc -u IP 9999}.

{\bf Watch using {\tt tcpdump -i eth1 -XXX} how the packets transferred from PC to Red Pitaya can be seen on the interface}

{\bf Share data encoded on more than 1-byte and observe the endianness issue of the byte organization between
sender, network and receiver}

{\bf Write a Makefile compatible with both cross-compilation and compilation (hint: use {\tt CC ?=})}.
  
{\bf Manually cross-compile PocketSDR from \url{https://github.com/tomojitakasu/PocketSDR} so it executes on the Red Pitaya.}
  
\subsection{Packaging for Buildroot}

Rather than cross-compiling manually, a software available on the internet might be compiled as
a Buildroot package. To do so
\footnote{\url{https://bootlin.com/~thomas/site/buildroot/adding-packages.html}}:

\begin{itemize}
\item add an entry in \verb~package/Config.in~ with your package directory, e.g. \verb~source "package/mypackage/Config.in"~
\item create under the \verb~package/~ directory a new \verb~mypackage~ directory that will hold the configuration files
\item create a \verb~Config.in~ file in \verb#package/mypackage# with the description of your package and the variable
associated with its selection:
\begin{verbatim}
config BR2_PACKAGE_MYPACKAGE
        bool "mypackage"
        help
          A new package example
\end{verbatim}
\item create a \verb~mypackage.mk~ file with the rules for compiling the package. For a simple {\tt makefile} based
compilation:
\begin{verbatim}
RPITX_VERSION = mypackage_hash
RPITX_SITE = $(call github,project,mypackage,$(RPITX_VERSION))

RPITX_INSTALL_STAGING = YES
RPITX_LICENSE = GPL-3.0
RPITX_LICENSE_FILES = GNU_General_Public_License.txt

define MYPACKAGE_BUILD_CMDS
    $(TARGET_MAKE_ENV) $(TARGET_CONFIGURE_OPTS) $(MAKE) -C $(@D)/src
endef

define MYPACKAGE_INSTALL_STAGING_CMDS
endef

define MYPACKAGE_INSTALL_TARGET_CMDS
        $(INSTALL) -D -m 0755 $(@D)/mypackage $(TARGET_DIR)/usr/bin
endef

$(eval $(generic-package))
\end{verbatim}
\end{itemize}

Notice that the {\tt Makefile} must be cross-compilation compatible, e.g. using {\tt CC ?=} rather than {\tt CC =} so
that an external definition of the variable {\tt CC} overrides the default compiler selection. The {\tt BUILD\_CMDS} includes
the {\tt \$(MAKE) -C dir} meaning that the {\tt make} configuration file is search for in the {\tt dir} subdirectory.

If the additional package is included in a new {\tt BR2\_EXTERNAL} repository, {\bf multiple external repositories}
can be used by defining

\verb~export BR2_EXTERNAL=$HOME/oscimp_br2_external/:$HOME/redpitaya~

Store library compilation output to the {\tt staging} directory (found in {\tt output/staging}) so the subsequent packages 
can link to its output. Notice that
\footnote{\url{https://buildroot.org/downloads/manual/manual.html}} ``staging/ is a symlink to the target toolchain sysroot inside host/, which exists for 
backwards compatibility'' so that the actual target sysroot is located in the host directory~! (see \\
{\tt output/host/arm-buildroot-linux-gnueabihf/sysroot})

{\bf Generate the packages needed for cross-compiling PocketSDR using Buildroot. Notice that PocketSDR is using two external git sub-modules
which is not readily supported by Buildroot which prefers generating separate packages for {\tt libfec} and {\tt libLDPC}
\footnote{\url{https://buildroot.org/downloads/manual/manual.html} states that ``LIBFOO\_GIT\_SUBMODULES can be set to YES to create 
an archive with the git submodules in the repository. This is only available for packages downloaded with git (i.e. when LIBFOO\_SITE\_METHOD=git). 
Note that we try not to use such git submodules when they contain bundled libraries, in which case we prefer to use those libraries from 
their own package''}. When working in the PocketSDR directories to update Makefiles in order to make them compatible with cross-compilation,
all changes can be saved and applied when running Buildroot compilation by executing \\ {\tt git diff > br2\_external/package/pocketsdr/001-makefiles.patch}: the {\tt 001-makefiles.patch} will be applied by Buildroot prior to running the compilation commands.}

\section{Hardware access using a C program}

Accessing hardware resources from the shell is functional, but 1/ is slow since it
must reach through all the abstraction layer between userspace and the hardware
including the kernel and 2/ assumes that someone before us has already implemented the 
driver to control the hardware. We will here discuss how to control hardware without
relying on the kernel -- an efficient means for quickly prototyping as we are used to
when working with microcontrollers not running an operating system but against the
principles when developing with GNU/Linux -- which will later be useful when qualifying
the latencies introduced by the operating system.

{\tt buildroot} has provided a consistent development framework and most significantly
a cross-compiler in {\tt output/host/bin} named {\tt arm-linux-gcc}. 
The datasheet describing the Zynq registers is found at \\
\url{https://docs.amd.com/r/en-US/ug585-zynq-7000-SoC-TRM/}.
As usual, we find there all the addresses of the registers used to interact with
the hardware as well as their usage. As usual, such a complex datasheet hides the details
that make the difference between a functional and a failing program. Here, the surprise
lies, as was the case with the STM32, on the {\bf need to activate the clock driving the
GPIO peripheral}: if not activated, reading or writing to a register associated with this
peripheral will lead to a null result (0x00). The solution is provided, for an MIO, at
\url{https://support.xilinx.com/s/question/0D52E00006hpU2mSAE/using-genericuio-with-ps-gpio}
% \url{https://forums.xilinx.com/t5/Embedded-Linux/Zynq-mmap-GPIO/td-p/368601} 
which explains
which clock to activate (bit 22 of the APER\_CLK\_CTRL register). Furthermore, remember
that a {\tt unsigned int *h;} pointer is incremented with 4-byte steps, so that
{\tt h+4} points towards the address of {\tt h} incremented by {\bf 16 bytes} and not 4~bytes
as would be erroneously expected. Thus, we find easier to define {\tt unsigned char *h;} which
avoids any surprise when incrementing addresses.

The example below allows for blinking the two LEDs connected to the two MIOs driven by
the processing system.

\begin{lstlisting}[language=C,basicstyle=\footnotesize\ttfamily]
// arm-linux-gcc -Wall -o led led.c

#include <stdio.h>
#include <fcntl.h>    // O_*
#include <unistd.h>   // close
#include <stdlib.h>   // atoi
#include <stdint.h>   // uint32_t
#include <sys/mman.h> // mmap

//#define TYPE int  // si on definit h comme int*, alors l'increment est en multiple de 4 !
#define TYPE char

int main(int argc, char** argv)
{ const int base_addr1= 0xF8000000;
  const int base_addr2= 0xe000a000;
  TYPE *h    = NULL; // int *h = NULL;
  int map_file        = 0;
  unsigned int page_addr,page_offset;
  unsigned page_size=sysconf(_SC_PAGESIZE);

  uint32_t value= 0x80; // 1 ou 128 pour LED orange ou rouge
  if (argc >= 2) value = atoi(argv[1]);
  printf("value = 0x%x\n",value);

  page_addr  = base_addr1 & (~(page_size-1));
  page_offset= base_addr1-page_addr;
  printf("page=0x%x size=0x%x offset=0x%x\n",page_addr,page_size,page_offset);
  
  // mmap the device into memory 
  map_file = open("/dev/mem", O_RDWR | O_SYNC); 
  h=mmap(NULL,page_size,PROT_READ|PROT_WRITE,MAP_SHARED,map_file,page_addr);
  if (h<0) {printf("mmap pointer: %x\n",(int)(h));return(-1);}

  // exemple baremetal pour debloquer acces SLCR registers
  // https://forums.xilinx.com/xlnx/attachments/xlnx/EDK/29780/1/helloworld.c
  printf("@      -> %x %x \n",(unsigned int)h,(unsigned int)(h+0x0c));
  printf("init   -> %x \n",(*(unsigned int*)(h+0x0c/sizeof(TYPE))));
  *(unsigned int*)(h+0x04/sizeof(TYPE))=0x767b; // lock
  printf("lock   -> %x\n",(*(unsigned int*)(h+0x0c/sizeof(TYPE))));
  *(unsigned int*)(h+0x08/sizeof(TYPE))=0xDF0D; // unlock
  printf("unlock -> %x\n",(*(unsigned int*)(h+0x0c/sizeof(TYPE))));
  
  // IL FAUT L'HORLOGE DE GPIO ! bit 22 de 0xF8000000+0x0000012C, sinon meme lecture echoue (0x00)
  // https://forums.xilinx.com/t5/Embedded-Linux/Zynq-mmap-GPIO/td-p/368601
  printf("ck -> %x\n",(*(unsigned int*)(h+0x12c/sizeof(TYPE))));
  *(unsigned int*)(h+0x12c/sizeof(TYPE))|=(1<<22);
  printf("ck -> %x\n",(*(unsigned int*)(h+0x12c/sizeof(TYPE))));

  page_addr = (base_addr2 & (~(page_size-1)));
  page_offset = base_addr2 - page_addr;
  printf("%x %x\n",page_addr,page_size);
  h=mmap(NULL,page_size,PROT_READ|PROT_WRITE,MAP_SHARED,map_file,page_addr);
  if (h<0) {printf("mmap pointer: %x\n",(int)(h));return(-1);}
  //write_reg(0xE000A000, 0x00000000, 0x7C020000); //MIO pin 9 value update
  //write_reg(0xE000A000, 0x00000204, 0x200);	 //set direction of MIO9
  //write_reg(0xE000A000, 0x00000208, 0x200);	 //output enable of MIO9
  //write_reg(0xE000A000, 0x00000040, 0x00000000); //output 0 on MIO9
  //data = read_reg(0xE000A000, 0x00000060);       //read data on MIO

  // https://docs.amd.com/r/en-US/ug585-zynq-7000-SoC-TRM/
  *(unsigned int*)(h+page_offset+0)=(value);     // GPIO programming sequence :p.386
  printf("hk: %x\n",*(unsigned int*)(h+page_offset));
  *(unsigned int*)(h+page_offset+0x204/sizeof(TYPE))=value;  // direction
  printf("hk+204: %x\n",*(unsigned int*)(h+page_offset+(0x204)/sizeof(TYPE)));
  *(unsigned int*)(h+page_offset+0x208/sizeof(TYPE))=value;  // output enable
  printf("hk+208: %x\n",*(unsigned int*)(h+page_offset+(0x208)/sizeof(TYPE)));
  *(unsigned int*)(h+page_offset+0x040/sizeof(TYPE))=value;  // output value
  printf("hk+40: %x\n",*(unsigned int*)(h+page_offset+0x40)/sizeof(TYPE));
  close(map_file);
  return 0;
}
\end{lstlisting}

This example will be used multiple times, whether for qualifying the latencies of
Xenomai or accessing the registers configured in the programmable logic (PL) of the
Zynq.

% /dev/xdevcfg :  cat system_wrapper.bit > /dev/xdevcfg 
% avec http://antonpotocnik.com/?p=487360
% comment faire desormais pour acceder au PL depuis un noyau mainline ?

Such direct memory accesses from userspace are possible -- without writing a C
program handling the {\tt /dev/mem} interface and convert real to virtual
addresses ({\tt mmap}) -- thanks to the {\tt devmem}
\footnote{sources available in {\tt buildroot} at {\tt output/build/busybox-1.30.1/miscutils/devmem.c}}
software. Hence for example, unlocking the Zynq configuration registers can be tested with
{\footnotesize
\begin{verbatim}
redpitaya> devmem 0xf8000004 32 0x767b   # password to lock
redpitaya> devmem 0xf800000c             # check status: 1=locked
0x00000001
redpitaya> devmem 0xf8000008 32 0xdf0d   # password to unlock
redpitaya> devmem 0xf800000c             # check status: 0=unlocked
0x00000000
\end{verbatim}
}

\subsection{Reconfiguring pin functions -- handling hardware interrupts}

The Red Pitaya only provides two GPIO controlled from the PS -- the MIOs -- to which LEDs are connected.
These pins are hence not usable as inputs or for any other use than visual information. Other
MIO pins are available, but their default configuration ({\tt ps7\_init.c} file generated by Vivado or provided 
with the Linux kernel for starting the system) is dedicated to alternative functions (SPI, UART).

We wish to use some of these GPIOs to illustrate hardware interrupt handling. The GPIO functionality
must hence be restored, for example when considering the SPI signal provided on MIO10, pin 3 of connector
E2.

The Xilinx datasheet is clear on this topic: the function of a pin is defined with 3 bits named
L0\_SEL, L1\_SEL and L2\_SEL. We find as many configuration registers, and hence such bits, than
pins to be configured. We identify the address associated with MIO10, check that its initial default
configuration is indeed SPI communication, and restore it as a GPIO:

{\footnotesize
\begin{verbatim}
redpitaya> cd /sys/class/gpio/
redpitaya> echo "916" > export 
redpitaya> echo "out" > direction 
redpitaya> devmem 0xF800012C # GPIO clock active
0x00500444
redpitaya> devmem 0xF8000728 # A = SPI
0x000016A0
redpitaya> devmem 0xF8000728 32 0x00001600
redpitaya> devmem 0xF8000728 # 0= GPIO
0x00001600
redpitaya> echo "1" > value 
redpitaya> echo "0" > value 
\end{verbatim}
}

\includegraphics[width=0.49\linewidth]{spi1.png}
\includegraphics[width=0.49\linewidth]{spi2.png}

Following these commands, the MIO10 pin (EP2 pin 3, Fig. \ref{pinout}) is indeed set as a GPIO, as
demonstrated with the following Linux kernel module that uses the kernel {\em timer} to 
alternate the pin state. 

\begin{figure}[h!tb]
\begin{center}
\includegraphics[width=0.49\linewidth]{Extension_connector.png}
\end{center}
\caption{Redpitaya extention port pinout, from {\tt https://redpitaya.readthedocs.io/en/latest/developerGuide/\\hardware/125-14/extent.html}. MIO10 is indicated here as SPI\_MOSI or pin 3 of E2.}
\label{pinout}
\end{figure}

Notice that the argument provided to {\tt gpio-lib} to configure 
the GPIO is the same than the one we provided earlier to {\tt /sys/class/gpio}, namely 
906+GPIO index (here 10) for pre-6 kernels and 512+GPIO index for 6.x kernels:

\lstinputlisting[language=C]{programs/4mymod_version_gpiolib.c}

Similarly, this pin can trigger an asynchronous event (interrupt) its its voltage
changes. An interrupt is handled in the sample kernel module as follows:

\lstinputlisting[language=C]{programs/5mymod_int.c}

which is further optimized with interrupt handling at the kernel space aimed at sending
a signal that might have registered with such a service:

\lstinputlisting[language=C]{programs/6mymod_int.c}

in order to later share the information to userspace:

\lstinputlisting[language=C]{programs/6user_prog.c}

% from ~/enseignement/ufr/M1_info/SeanceA_kernel/kernel_module_Red Pitaya.tex

% 3. verifier ioremap
% verifier ioctl avec 64 bits

% # cat /proc/iomem 
% ...
% 01c20800-01c20bff : /soc@01c00000/pinctrl@01c20800
% 01c20c90-01c20c9f : /soc@01c00000/watchdog@01c20c90

\section{Kernel programming: modules and communicating with userspace}

Our objective in this section is to understand how to communicate between userspace
and hardware by benefitting from kernel features. Although we have seen that various
hardware abstractions ({\tt script.bin} or {\tt devicetree.dtb}) aim at hiding
the hardware details to the developer, we will he consider the electronics engineer
perspective who understands hardware and wishes to control from the Linux kernel
to provide standardized interfaces to the user.

\subsection{Why work at the kernel level?}

Considering the abstraction layers between a userspace software and hardware, the
kernel aims at keeping consistent the various requests to control the few hardware resources
from the many possible applications. The kernel makes sure each task is
given some time to execute (scheduler), that a unique task can access a hardware
resource at any given time, and that messages read from the hardware are communicated
to all userspace applications which requested access.

A module \cite{ldd3} is a piece of software that connects to the kernel
to add features. Initially introduced with Linux kernel version 1.2 (March 1995),
it avoids compiling the whole kernel when adding new features such as support
for a new hardware peripheral. Most hardware peripherals could be accessible
from userspace (unless Direct Memory Access or interrupts are involved) but doing
so would waste the consistency of a single software driving a peripheral provided
by the kernel supervision (for example by preventing multiple copies of a module
accessing the same hardware being loaded). Some features, such as the mentioned
Direct Memory Access or interrupts, are just not accessible outside the kernel.

When configuring the kernel, selecting a module is achieved by setting to
{\tt M} such options when running {\tt make menuconfig}. Compiling modules
is completed with {\tt make modules} followed by {\tt make modules\_install} 
which copies the compilation result to {\tt /var/lib/modules}.

These options are accessed when using buildroot by executing {\tt make linux-menuconfig}
so that Buildroot's {\tt make} will compile both kernel and modules, and store the
result in the image to be stored on the SD card.

Any driver developer must understand the architecture of a kernel module in order to
provide the interfaces expected by a user willing to access hardware without understanding
the details of hardware configuration.

The two classes of kernel modules are designed for transferring data blocks (for example
to a mass storage interface such as a hard disk), or as individual bytes
\footnote{\url{https://appusajeev.wordpress.com/2011/06/18/writing-a-linux-character-device-driver/}} 
We will focus on the latter, easier to grasp and meeting most sensor applications while
handling hardware peripherals easier to run during demonstrations.

All communication is achieved through pseudo-files located in the {\tt /dev/} directory.
A file can be opened, closed, and data can be written to or read from the file. An additional
method which does not have any meaning in terms of file access is the input output
configuration: {\tt ioctl} or (Input/Output Control).


\subsection{Structure of a kernel module and compilation}

The minimum needed for linking a module to a Linux kernel are the functions for initializing
resources and a function for freeing these resources. The following examples provides a 
skeleton and become familiar with the Linux kernel tree structure on the one hand, and compilation
method through a {\tt Makefile} requesting methods associated with the kernel being run on the target
platform.

A minimalistic kernel module is provided below, for becoming familiar with the initialization
function and compilation mechanism as well as linking the module with the kernel.

\begin{lstlisting}[language=C,label=code1,caption={First example of kernel module}]
#include <linux/module.h>       /* Needed by all modules */
#include <linux/kernel.h>       /* Needed for KERN_INFO */
#include <linux/init.h>         /* Needed for the macros */

static int __init hello_start(void){printk(KERN_INFO "Hello\n");return 0;}
static void __exit hello_end(void) {printk(KERN_INFO "Goodbye\n");}

module_init(hello_start);
module_exit(hello_end);

MODULE_LICENSE("GPL");
\end{lstlisting}

The kernel module defines the entry and exit functions. When a module is linked to the kernel
({\tt insmod my\_mod.ko}), the {\tt init()} method is called. When the module is unlinked from
the kernel ({\tt rmmod my\_mod}), the {\tt exit} method is called. Both methods are macros
calling the {\tt init\_module()} and {\tt cleanup\_module()} functions as defined at
{\tt linux/init.h} in the Linux kernel source code.

\noindent{\fbox{\parbox{\linewidth}{While historically the {\tt .o} object was directly linked
to the kernel, since its 2.6 version the Linux kernel requires additional informations to link
with a {\em kernel object} which differs between the {\tt .o} and {\tt .ko}. Although the Makefile
declares a dependency on the {\tt .o} output, the {\tt .ko} file is indeed the one to be loaded
using {\tt insmod} as described at \url{https://tldp.org/HOWTO/Module-HOWTO/linuxversions.html}: 
``In Linux 2.6, the kernel does the linking. A user space program passes the 
contents of the ELF object file directly to the kernel. For this to work, the ELF object image 
must contain additional information. To identify this particular kind of ELF object file, 
we name the file with suffix ".ko" ("kernel object") instead of ".o" For example, the serial 
device driver that in Linux 2.4 lived in the file serial.o in Linux 2.6 lives in the file 
serial.ko.''}}

A module is necessarily linked with a given and known version of the kernel, and will have
to be recompiled whenever the latter is updated. Buildroot provides a consistent development
framework including the cross-compilation toolchain and the kernel sources, which are found at
{\tt output/build/linux-}{\em version} of buildroot. In our case, the Linux kernel is fetched
using {\tt git} and the directory name is postfixed with the hash key.

As a first step, we make sure the compiler is in the executable path:
\begin{verbatim}
export PATH=$PATH:$HOME/.../buildroot/output/host/usr/bin/
\end{verbatim}

\noindent
then we compile the module -- assumed to be named {\tt mymod.c}
\footnote{{\bf do not call the module} {\tt kernel.c}: doing so will lead
to a successful compilation, but by using an object that does not result
from our own source code but from a file located in the kernel tree structure!}
to become {\tt mymod.ko} -- using the following {\tt Makefile}:

\begin{lstlisting}
obj-m	+= mymod.o

all:
	make ARCH=arm CROSS_COMPILE=arm-linux- -C \
  $(BR)/output/build/linux-[...] M=$(PWD) modules

clean:
	rm mymod.o mymod.ko
\end{lstlisting}

This {\tt Makefile} is analyzed as follows:
\begin{enumerate}
\item it calls itself another Makefile since executing the {\tt make} command
\item this other Makefile is located in the sources of the kernel we will be linking and
is selected thank to the {\tt -C directory} option which indicates the location of the 
targeted configuration,
\item we request the {\tt modules} method since this command could be summarized as
{\tt make modules} with a few additional options...
\item ... amongst which the cross-compilation towards and ARM target, hence the {\tt ARCH}
option and the compiler prefix in {\tt CROSS\_COMPILE}.
\item Finally, compiling a module requires providing the location in which the sources are
located: here we are considering the working directory so that 
\verb~M=$(PWD)~ in this case.
\item If the same module is to be compiled on the PC (host architecture), we remove
{\tt ARCH} and {\tt CROSS\_COMPILE} to only keep {\tt make -C kernel\_directory M=\$(PWD) modules} (see below)
\end{enumerate}

% 200331
%La seule m\'ethode importante est \'evidemment {\tt all}, qui fournit en argument de {\tt make}
%le pr\'efixe du compilateur (ici {\tt arm-buildroot-linux-uclibcgnueabihf-} puisque le compilateur
%g\'en\'er\'e par {\tt buildroot} se nomme \\
%{\tt arm-buildroot-linux-uclibcgnueabihf-gcc}), l'emplacement
%des sources du noyau, et l'architecture cible. 
%
%Ce makefile est un peu complique' car il s'appuie sur le Makefile des sources de
%Linux qu'on appelle avec l'argument -C.
%
%Ce que je raconte quand je fais cours : il faut lire ce Makefile de la facon suivante
%* make -C repertoire signifie que je vais lancer le Makefile qui se trouve dans repertoire,
%  ici les sources du noyau linux
%* quelle methode de ce Makefile je lance ? la methode modules car on voit qu'on fait un
%  make modules
%donc pour un module que tu fabriques pour le noyau de ton PC, ca se resume a 
%make -C /usr/src/linux modules
%* ici on cross-compile pour cible ARM, donc on ajoute les deux argument ARCH= et CROSS_COMPILE=
%* finalement, le Makefile de Linux veut savoir ou` se trouvent les sources du noyau en cours de 
%compilation : c'est le repertoire courant donc M=$(PWD)
\noindent{\fbox{{\parbox{\linewidth}{A kernel module is not only compiled to the embedded target
board, but also possibly to the host as long as low-level hardware is not accessed. In order
to know the kernel version running on the PC, run the {\tt uname -a} command. Make sure the kernel
sources are available, or at least the associated headers and configuration, for example with
the package {\tt linux-headers-amd64} when using Debian GNU/Linux. On the PC, the Makefile will look like
\\
{\tt
obj-m += t.o\\
all:\\
\hspace*{1cm} make -C /usr/src/linux-headers-6.10.4-amd64 M=\$(PWD) modules
}
}}}
making sure {\bf not} to use the {\tt common} extension which would not provide the 
configuration of this particular kernel for the current host architecture.

We observe how powerful Buildroot is in providing a consistent working environment with the
cross-compilation toolchain ({\tt CROSS\_COMPILE=arm-buildroot-linux-uclibcgnueabihf-)} and the
kernel tree structure \\
({\tt buildroot-2025.05.1.tar.gz/output/build/linux-xilinx-v2024.1})
which is used when compiling the module. The principle of the {\tt Makefile} is to add the list
of kernel modules under development ({\tt mymod}) to the list of kernel modules, and run the
compilation of the kernel with {\tt make}. The other kernel objects having been already compiled, only
our current module will be generated.

\noindent\fbox{\parbox{\linewidth}
{{\bf 
\includegraphics[width=.7cm]{danger} %\danger 
Warning}: when using a board running a kernel other than the one provided by Buildroot, synchronizing
the kernel version with the one running on the board is mandatory. Doing so is achieved by fetching
on the embedded board the Linux configuration found in {\tt /proc/configs.gz}, copy this file in the
{\tt .config} configuration file of the kernel being compiled by Buildroot, run {\tt make linux-menuconfig}
to tell Buildroot of the update, and finally {\tt make} from the top Buildroot location to recompile
the full kernel and dependencies. Once the compilation is completed, the module we will compile from now
on will be compatible with the kernel running on the board.}}

The module is linked to the kernel with {\tt insmod mymod.ko} and removed with {\tt rmmod mymod}. The list
of modules linked to the kernel is displayed with {\tt lsmod}. The messages from the kernel, including its modules,
are displayed on the console (in the case of the Red Pitaya, on the serial port) or in the kernel logs displayed
with {\tt dmesg} or {\tt cat /var/log/messages} (or {\tt /var/log/syslog}).

The result of the kernel module compilation is transferred to the Red Pitaya board (NFS or {\tt scp})
to be linked to the kernel with {\tt insmod mymod.ko}. We validate that the command was
acknowledged by displaying the log messages with {\tt dmesg}:

\begin{verbatim}
# insmod mymod.ko 
# dmesg | tail -1
[  249.879087] Hello
# rmmod mymod.ko 
# dmesg | tail -1
[  256.571319] Goodbye
\end{verbatim}

These information are also found in {\tt /var/log/messages} which can be constantly monitored using
{\tt tail -f}.

% This example is derived from the kernel sources as found at {\tt modules/example/example.c}.

When a module compilation fails, it is fundamental to remember that a kernel module is linked against
a given kernel configuration. If this configuration is lost, is can always be recovered with
{\tt scripts/extract-ikconfig} whose output should be saved in a {\tt .config} file allowing to
regenerate a working environment consistent with the one expected by the kernel being executed.

The communication capabilities of our module are however limited: we shall extend such capabilities
with a character device interface ({\em char device}).

\subsection{Communication through {\tt /dev}: interacting with the user}

A module only able to load and unload is hardly usable. In order to interact with the user, we must
provide methods for reading or writing. The standard location of the pseudo-files allowing such
interaction between the kernel and userspace is in the {\tt /dev} directory. These files act as
pipes between the two spaces (kernel and user) through which single bytes (characters) are shared 
(as opposed to data blocks). The administrator (root or {\tt sudo} on the host PC) can create such 
pseudo-files using the command {\tt mknod /dev/jmf c 90 0} with the arguments indicating that we are creating 
a {\em character device} ({\tt c}) identified with the major number 90 (kind of peripheral) and minor number 0 
(index in the list of such peripherals). On the PC, adding {\tt -m 666} to the node creation will allow for
read/write permissions to all users: {\tt mknod /dev/jmf c 90 0 -m 666}.

\noindent\fbox{\parbox{\linewidth}{
\includegraphics[width=.7cm]{danger} %\danger 
In case a resource busy message is displayed when inserting the module (error -16 -- EBUSY), the major
number 90 might alread be used by another module: even if it does not appear in the list of 
entries in {\tt /dev}, the list of peripherals in {\tt /proc/devices} will provide all the major numbers
already allocated to kernel modules. For example on the Red Pitaya, major number 90 is allocated to the
{\tt mtd} module. On a host PC, the 99 major number is allocated to {\tt ppdev}.}}

The structure defining which function is called when a system call is transferred from userspace to kernel
space (function pointer) when opening or closing the pipe ({\tt open} and {\tt close} in userspace) as well
as reading and writing data through the pipe ({\tt write} and {\tt read} system calls from userspace
respectively) is name {\tt file\_operations}.

In the example below, writing to {\tt /dev/jmf} will lead to a message in the kernel logs including the string
that was written by the user:

\begin{lstlisting}[language=C]
// mknod /dev/jmf c 90 0
#include <linux/module.h>       /* Needed by all modules */
#include <linux/kernel.h>       /* Needed for KERN_INFO */
#include <linux/init.h>         /* Needed for the macros */
#include <linux/fs.h>           // define fops
#include <linux/uaccess.h>

static int dev_open(struct inode *inod,struct file *fil);
static ssize_t dev_read(struct file *fil,char *buff,size_t len,loff_t *off);
static ssize_t dev_write(struct file *fil,const char *buff,size_t len,loff_t *off);
static int dev_rls(struct inode *inod,struct file *fil);

char buf[15]="Hello read\n\0";

int hello_start(void); // declaration pour eviter les warnings
void hello_end(void);

static struct file_operations fops=
{.read=dev_read,
 .open=dev_open,
 .write=dev_write,
 .release=dev_rls,
};

int hello_start()  // init_module(void) 
{int t=register_chrdev(90,"jmf",&fops);  // major = 90
 if (t<0) printk(KERN_ALERT "registration failed\n");
    else printk(KERN_ALERT "registration success\n");
 printk(KERN_INFO "Hello\n");
 return t;
}

void hello_end() // cleanup_module(void)
{printk(KERN_INFO "Goodbye\n");
 unregister_chrdev(90,"jmf");
}

static int dev_rls(struct inode *inod,struct file *fil)
{printk(KERN_ALERT "bye\n");return 0;}

static int dev_open(struct inode *inod,struct file *fil)
{printk(KERN_ALERT "open\n");return 0;}

static ssize_t dev_read(struct file *fil,char *buff,size_t len,loff_t *off)
{int readPos=0;
 printk(KERN_ALERT "read\n");
 while (len && (buf[readPos]!=0)) 
   {put_user(buf[readPos],buff++);
    readPos++;
    len--;
   }
 *off += readPos; // avoid repeating the output when cat < /dev/XXX
 return readPos;
}

static ssize_t dev_write(struct file *fil,const char *buff,size_t len,loff_t *off)
{int m=0,mylen;
 printk(KERN_ALERT "write ");
 if (len>14) mylen=14; else mylen=len;
 for (m=0;m<mylen;m++) get_user(buf[m],buff+m);
 // DO NOT try with vvv: kernel panic since accessing usespace from kernel space
 // for (m=0;m<mylen;m++) buf[m]=buff[m];
 buf[mylen]=0;
 printk(KERN_ALERT "%s \n",buf);
 return len;
}

module_init(hello_start);
module_exit(hello_end);
\end{lstlisting}

This module is loaded and its functions called using
\begin{multicols}{2}
{\footnotesize
\begin{verbatim}
# echo "Hello" > /dev/jmf 
[  258.408111] open
[  258.414953] write 
[  258.416874] Hello
[  258.416874]  
[  258.420911] bye
# dd if=/dev/jmf bs=10 count=1
[  282.774521] open
[  282.776557] read
Hello[  282.778454] bye

0+1 records in
0+1 records out
# echo "toto" > /dev/jmf [  290.922228] open

[  290.924908] write 
[  290.926933] toto
[  290.926933]  
[  290.930904] bye
# dd if=/dev/jmf bs=10 count=1
[  293.484720] open
[  293.486761] read
toto[  293.488659] bye

0+1 records in
0+1 records out
\end{verbatim}
}
\end{multicols}

\noindent{\bf This example includes a modification to the module shown previously so that
the displayed message is the one read previously: implement this modification.}

This module implements communication between user space and kernel space through the pseudo-file
named {\tt /dev/jmf}. A file located in the {\tt /dev} is not a file for storing data but a link
allowing for sharing data and system calls between user space and kernel space. This file was created
with {\tt mknod /dev/jmf c 90 0} with 90 its major number identified and 0 its minor number identifier.
The major number remains constant for all files related to a given type of peripheral, while the
minor number indexes which peripheral is considered by being incremented for each new instance.
We check that the association of the communication peripheral with the module is correct by reading
{\tt /proc/devices} which includes the {\tt jmf} entry associated with the major number we selected.
(90, which is free as verified with {\tt ls -l /dev/}). However, the {\tt /dev/} is often dynamically
generated during the boot sequence, and any new entry will have to be recreated after a reboot sequence.

An alternate solution to explicitly providing the communication entry in {\tt /dev} by manually selecting
the major number is to use the {\em misc device} class. Doing so is achieved by creating the global variable
\begin{lstlisting}[language=C]
 struct miscdevice jmfdev;
\end{lstlisting}
and in the init function defining:
\begin{lstlisting}[language=C]
 jmfdev.name = "jmf";  // /dev/jmf using the major number of the misc class
 jmfdev.minor = MISC_DYNAMIC_MINOR;
 jmfdev.fops = &fops;
 jmfdev.mode = 0666; // permissions for the user to access
 misc_register(&jmfdev);  // dynamic creation of the /dev/jmf entry
\end{lstlisting}
which is concluded in the exit function with
\begin{lstlisting}[language=C]
 misc_deregister(&jmfdev);
\end{lstlisting}

\subsection{The {\tt ioctl} system call}

After defining the systems calls matching the requirement that all peripherals must
appear as files, it occurred that one more system call had to be added that breaks this
initial Unix design philosophy: this new system call is called {\tt ioctl} for input/output
control. This system call has been added to complement functionalities that cannot be solved
with just {\tt read} and {\tt write}, namely configuration of the peripherals (for example,
tuning the sampling rate or the data size of a sound card).

The previous example listing can hence be complemented with the following functions:

\begin{lstlisting}[language=C]
// static int dev_ioctl(struct inode *inod,struct file *fil, unsigned int cmd, 
//      unsigned long arg); // pre-2.6.35
static long dev_ioctl(struct file *f, unsigned int cmd, unsigned long arg); // actuel

static struct file_operations fops=
{.read=dev_read,
 .open=dev_open,
 .unlocked_ioctl=dev_ioctl,
 .write=dev_write,
 .release=dev_rls,};

// static int dev_ioctl(struct inode *inod,struct file *fil, unsigned int cmd, 
//    unsigned long arg) // pre-2.6.35
static long dev_ioctl(struct file *f, unsigned int cmd, unsigned long arg)
{switch ( cmd ) 
 {case 0: printk(KERN_ALERT "ioctl0");break;
  case 1: printk(KERN_ALERT "ioctl1");break;
  default:printk(KERN_ALERT "unknown ioctl");break;
 }
 return 0;
}
\end{lstlisting}

Executing the following program -- we are not aware of any way of calling {\tt ioctl()} from the 
shell without writing a dedicated C program -- which we call {\tt ioct}

\begin{lstlisting}[language=C]
#include <fcntl.h>      /* open */ 
#include <unistd.h>     /* exit */
#include <sys/ioctl.h>  /* ioctl */
#include <stdio.h>
#include <stdlib.h>

#define IOCTL_SET_MSG 0
#define IOCTL_GET_MSG 1
#define DEVICE_FILE_NAME "/dev/jmf"

ioctl_set_msg(int file_desc, char *message)
{ int ret_val;
  ret_val = ioctl(file_desc, IOCTL_SET_MSG, message);
  printf("set_msg message: %s\n", message);
}

ioctl_get_msg(int file_desc,char *message)
{ int ret_val;
  ret_val = ioctl(file_desc, IOCTL_GET_MSG, message);
  printf("get_msg message: %s\n", message);
}

main(int argc,char **argv)
{ int file_desc, ret_val;
  char msg[30] = "Message passed by ioctl\n";

  file_desc = open("/dev/jmf", 0);
  printf("%d\n",file_desc);
  if (argc>1) sprintf(msg,argv[1]);
  msg[4]=0;
  ioctl_set_msg(file_desc,msg);

  sprintf(msg,"Hello World");
  ioctl_get_msg(file_desc,msg);
  close(file_desc); 
}
\end{lstlisting}

will generate the following sequence:

\begin{verbatim}
# lsmod
Module                  Size  Used by
g_ether                44895  0 
# insmod mymod.ko
# ./ioctl 
# tail /var/log/messages 
Jan  1 00:11:45 buildroot kern.alert kernel: [  705.740703] registration success
Jan  1 00:11:45 buildroot kern.info kernel: [  705.745730] Hello
Jan  1 00:11:50 buildroot kern.alert kernel: [  710.805798] open
Jan  1 00:11:50 buildroot kern.alert kernel: [  710.808491] ioctl1
Jan  1 00:11:50 buildroot kern.alert kernel: [  710.818059] ioctl0
Jan  1 00:11:50 buildroot kern.alert kernel: [  710.820838] bye
\end{verbatim}

From a user perspective, the {\tt ioctl()} function is called exactly as 
{\tt read()} or {\tt write()} would be, but will handle 3~arguments: the file descriptor, the
ioctl index (consistent with the declaration of the kernel module, here 0 or 1) and possibly
an argument. The second page of the {\tt ioctl} manual is clear about these aspects. From the
kernel side, transferring arguments \footnote{\url{http://www.makelinux.net/ldd3/chp-6-sect-1}} 
towards userspace memory addresses requires using {\tt put\_user()} for a single scalar value
or {\tt copy\_to\_user} for an array (memory area from kernel space to user space).

\subsection{From virtual addresses to physical addresses}

Although many embedded systems do not bother with hardware handling of memory through a
Memory Management Unit -- MMU -- a multitasking operating system such as Linux benefits
significantly from the use of such a peripheral and justifies the selection of a microcontroller
powerful enough to be fitted with an MMU. Memory organization, handled by the MMU, should not
be cared of by the developer as long as memory allocation and access remain consistent,
{\em except} when accessing a known physical memory location, as for example required when 
reaching hardware configuration registers. In that case, we must perform the reverse conversion,
from the virtual memory representation to the real physical addressing. From a userspace perspective,
such a functionality is provided by {\tt mmap()}. From the kernel space, it is {\tt ioremap()}.

Electronics engineers are only aware of physical addresses for reaching a peripheral: this is
the address associated with each register in the microcontroller datasheet. The computer
scientist only knows virtual addresses, as handled by the operating system running on top of the
MMU. We must provide a link between these two worlds in order to let the communicate.
% Previous paragraph : to let the ??? communicate ?

At the kernel level, the instruction \cite[chap.9]{ldd3} in charge of this link is
{\tt void *ioremap(unsigned long phys\_addr, unsigned long size);}

We will illustrate its use by complementing the methods called when inserting a module
to switch on a LED, and switch off the LED when the module is removed from memory. This example
is for demonstration purposes only since in a practical application, the module
{\tt linux-*/drivers/gpio/gpio-xilinx.c} implements such functionalities, hiding the subtleties
of hardware access from the developer.

The resources needed to understand the organization of the registers needed to configure
GPIOs of the 14-bit Red Pitaya are described at
\url{http://redpitaya.readthedocs.io/en/latest/developerGuide/125-14/extent.html}
with the help of the unavoidable -- but rather lengthy -- 
% \url{https://www.olimex.com/wiki/Bare_Metal_programming_A13#Do_something_useful.2C_Modify_the_hello_world_to_set_GPIO}
\url{https://docs.amd.com/r/en-US/ug585-zynq-7000-SoC-TRM/}.
The addresses of the registers related to GPIO configuration are described in the chapter 14
of the user manual.
% du processeur A13 disponible \`a 
% \url{http://free-electrons.com/~maxime/pub/datasheet/A13%20user%20manual%20v1.2%2020130108.pdf} (chapitre 
% 33). On y constate notamment que 36 registres sont assign\'es \`a la configuration de chaque port.

We complement the kernel module written so far with
\begin{lstlisting}[language=C]
static void __iomem *jmf_gpio;   //int jmf_gpio;
#define IO_BASE2 0xe000a000

int hello_start()  // init_module(void) 
{int delay = 1;
 unsigned int stat;
[... gestion horloge GPIO ...]

 if (request_mem_region(IO_BASE2,0x2E4,"GPIO test")==NULL)
     printk(KERN_ALERT "mem request failed");
 jmf_gpio = (void __iomem*)ioremap(IO_BASE2, 0x2E4);
 writel(1<<mio,jmf_gpio+0x204);
 writel(1<<mio,jmf_gpio+0x208);
 writel(1<<mio,jmf_gpio+0x40);

 return 0;
}

void hello_end() // cleanup_module(void)
{printk(KERN_INFO "Goodbye\n");
 release_mem_region(IO_BASE2, 0x2e4);
}
\end{lstlisting}

\noindent
after remembering to include the header files {\tt linux/io.h} and {\tt linux/ioport.h}.
The code can only be understood by reading \cite[p.1347]{zynq} partly reproduced here, with
the description of the various registers handling GPIO behaviour, and \cite[p.1580]{zynq} 
for the clock handling part.

%Tout d'abord, l'adresse de d\'epart de tous les registres, et le nombre d'octets allou\'es
%pour la configuration de chaque registre est fourni dans
%\begin{center}
%\includegraphics[width=0.49\linewidth]{a10_1.jpg}
%\end{center}
%
%Nous y trouvons l'adresse de base 0x01C20800 et l'assignation de 0x24=36 registres (octets) par 
%port.
%
%Ensuite, nous devons configurer le bit du port qui nous int\'eresse -- celui li\'e \`a une diode verte
%qui s'appelle PG9 -- en sortie. Le registre de configuration de direction du port est d\'ecrit ci-dessous
%(deux pages adjacentes de la documentation)~:
%
%\begin{center}
%\includegraphics[width=0.49\linewidth]{a10_2.jpg}
%\includegraphics[width=0.49\linewidth]{a10_3.jpg}
%\end{center}
%
%Nous y voyons 3 bits assign\'es \`a PG9, les bits 4 \`a 6. On v\'erifiera que la configuration 
%{\tt writel(0x10,jmf\_gpio+36*6+0x04);} met bien le bit 4 de ce registre \`a 1, configurant la broche
%en sortie. Noter que l'offset 0xDC=220 est bien \'egal \`a $36\times 6+4$.
%
%Finalement, la valeur plac\'ee sur le port est d\'efinie dans
%
%\begin{center}
%\includegraphics[width=0.49\linewidth]{a10_4.jpg}
%\end{center}
%
%Noter que l'offset 0xE8=232 vaut bien $36\times 6+0x10$. PG propose 12 bits de GPIO, d'o\`u la taille
%de ce registre, dont le 9\`eme \'el\'ement est bien index\'e en 0x200. La ligne r\'esultante pour
%allumer la LED est donc {\tt writel(0x0200,jmf\_gpio+36*6+0x10);}
%
%\noindent
%{\bf Observer la diode verte \`a cot\'e du connecteur audio lorsque le module est charg\'e et retir\'e
%de la m\'emoire.}

\subsection{Adding a periodic timer}

The Linux kernel provides access to timers for clocking tasks. We use such resources to make the
LED periodically blink and display a message in the kernel log files.

In order to inform the user that the Red Pitaya board is alive, we wish to display a heartbeat
pattern on the LED blinking periodically. The Linux kernel provides the timer as described at
\cite[chap.7]{ldd3} and illustrated in the following code: 

\begin{lstlisting}[language=C]
struct timer_list exp_timer;

static void do_something(struct timer_list *t)
{printk(KERN_ALERT "plop");
 mod_timer(&exp_timer, jiffies + HZ);
}

int hello_start()  // init_module(void) 
{...
 // init_timer_on_stack(&exp_timer); 
 // exp_timer.function = do_something;
 // exp_timer.data = 0;
 timer_setup(&exp_timer,do_something,0);
 exp_timer.expires = jiffies + delay * HZ; // HZ specifies number of clock ticks/second
 add_timer(&exp_timer);
...
}

void hello_end() // cleanup_module(void)
{...
 del_timer(&exp_timer);
}
\end{lstlisting}

\noindent with {\tt jiffies} the internal representation in the kernel of the current time.

We observe in \verb!dmesg | tail! that the timer is indeed periodically triggered, once every
second, as requested by \verb~mod_timer(&exp_timer, jiffies + HZ);~:

\begin{verbatim}
[ 5697.089643] registration success
[ 5697.093780] Hello
[ 5698.094083] plop
[ 5699.094146] plop
[ 5700.094227] plop
[ 5701.094286] plop
[ 5702.094348] plop
\end{verbatim}

%Userspace mptr = mmap(0, size, PROT_READ | PROT_WRITE, MAP_FILE | MAP_SHARED, fd, 0);

\subsection{Semaphore, mutex and spinlock \cite{blaess}}

A driver providing a {\tt read} method continually communicates with the userspace process
requesting the information (\verb~cat < /dev/mydriver~). Such a condition is not representative
of a practical application, in which the producer requires some time to generate data that
are the be sent to userspace: the {\tt read} method hence behaves as a data consumer. In
this producer-consumer architecture, some mechanism for blocking on the one hand the
call to {\tt read} as long as the data have not been produced, and on the other hand to make
sure the access to the memory area shared by the producer and consumer remains consistent, is needed.
The first functionality is provided by the semaphore, and the second by the mutex and its 
faster implementation (but using more resources) the spinlock.

\subsubsection{Semaphore}

The semaphore acts here as a counter remembering how many times data were produced. A consumer
trying to fetch data decrements the semaphore: if the semaphore is already 0, then the consumer
(here the {\tt read} method) is blocked until the producer has incremented the semaphore. A
semaphore incremented multiple times allows the consumer to obtain multiple datasets.

The definition of the constants and prototypes used by semaphores are summarized with
\begin{lstlisting}[language=C]
#include <linux/semaphore.h>                            
struct semaphore mysem;                               
sema_init(&mysem, 0);       // init the semaphore as empty 
down (&mysem); 
up (&mysem); 
\end{lstlisting}

{\bf Simulate the periodic production of datasets leading to the increment of the semaphore in
a timer handler. Block the call to {\tt read} as long as datasets have not been produced.}

{\bf What happens when reading from a peripheral accessible in the {\tt /dev/} entry when the
semaphore has already incremented multiple times?}

{\bf Amongst the five core functions of a communicating module -- init, exit, open, read, release -- 
which one allows producing datasets only when needed by the user? Demonstrate how the previous
program is possibly modified to produce datasets (incrementing the semaphore) only when requested
by the user, and not as early as loading the module.}
% open car avant open il n'y a pas de transmission
% => ne declencher l'acquisition que dans open, avec un flag dans open qui definit si on incremente le semaphore

\subsubsection{Mutex and spinlock}

Data produced are stored in a buffer memory for sharing information between producer and 
consumer. A mechanism must handle the consistency between writing by the producer and reading
by the consumer: the buffer memory must not be overwritten while being read, and its content
should not be read as it is being written. The MutEx ({\em MUTually EXclusive}) provides a binary
lock -- locked or unlocked -- preventing two processes from simultaneously accessing the same
memory area or variable. Before reaching a buffer memory, each process (producer or consumer)
locks the mutex, preventing the other process which might wish to lock the mutex itself from 
continuing execution. Once the operation -- reading or writing -- completex, the mutex is unlocked
and allows resuming of the waiting process.

The mutex mechanism allows for putting a task in a sleep state until the mutex is unlocked. Making
the task sleep frees some processor resources, but requires changing context which might be 
lengthy and requiring significant processing resources if the read and write operations should be
perform with short latencies. An alternative solution is the spinlock, which operates similarly but
keeps the task in an active state by continuously probing the state of the lock: in this case,
the processor is continuously active, but the latency is lowered and no context change is needed.

The definitions of constants and prototypes of mutexes are as summarized with
\begin{lstlisting}[language=C]
#include <linux/mutex.h>                            
struct mutex mymutex;
mutex_init(&mymutex);
mutex_lock(&mymutex);
mutex_unlock(&mymutex);
\end{lstlisting}

The definitions of constants and prototypes of spinlocks are as summarized with
\begin{lstlisting}[language=C]
#include <linux/spinlock.h>                            
static DEFINE_SPINLOCK(myspin);
spin_lock_init(&myspin);
spin_lock(&myspin);
spin_unlock(&myspin);
\end{lstlisting}

{\bf Demonstrate an implementation of a mutex designed to provide consistency between
reading and writing a common variable. The text displayed by a call to the {\tt read} function
should be the one stored by a call to {\tt write} with the addition of the value of a counter
periodically incremented to simulate the production of data.}

{\bf Instead of updating the content of an array of characters to be displayed in the timer
handler itself, create a task which is scheduled when time allows and is triggered by the timer:
place the processing steps requiring most computational power (mutex locking and unlocking,
updating the buffer content) in this task.}

\subsection{{\tt sysfs}}

The interaction between a module and a user is from now on no longer handled by a
pseudo-file (char or block device) in the {\tt /dev} directory but by a {\tt /sysfs} in
example \ref{sim3}: a directory in {\tt /sys} is dynamically created when loading the module.
Hence, a {\tt platform\_device} initializes a communication entry in {\tt /sys}, here called 
{\tt gpio-simple}. The {\tt probe} method is called upon initialization of the platform
device with the same driver identifier defined in the ``.name'' structure: the first argument
of {\tt platform\_device\_register\_simple} must be the same as the driver identifier:

\begin{lstlisting}[label=sim3,language=C]
#include <linux/module.h>
#include <linux/gpio.h>
#include <linux/platform_device.h>
#include <linux/err.h>

static struct platform_device *pdev;

static int gpio_simple_probe(struct platform_device *pdev)
{ 
  return 0;
}

static void gpio_simple_remove(struct platform_device *pdev)
{ 
}

static struct platform_driver gpio_simple_driver = {
        .probe          = gpio_simple_probe,
        .remove         = gpio_simple_remove,
        .driver = {
                .name   = "gpio-simple",
        },
};

static int __init gpio_simple_init(void)
{ int ret;
  ret = platform_driver_register(&gpio_simple_driver);
  if (ret) return ret;
  pdev = platform_device_register_simple("gpio-simple", 0, NULL, 0);
  if (IS_ERR(pdev))
    {platform_driver_unregister(&gpio_simple_driver);
     return PTR_ERR(pdev);
    }
 return 0;
}

static void __exit gpio_simple_exit(void)
{ 
  platform_device_unregister(pdev);
  platform_driver_unregister(&gpio_simple_driver);
}

module_init(gpio_simple_init)
module_exit(gpio_simple_exit)

MODULE_DESCRIPTION("GPIO simple driver");
MODULE_LICENSE("GPL");
\end{lstlisting}

\includegraphics[width=.7cm]{danger}Notice that the prototype of the {\tt .remove} function
was changed
\footnote{\url{https://github.com/torvalds/linux/commit/0edb555a65d1ef047a9805051c36922b52a38a9d}} around revision 6.10 of the Linux kernel: older versions returned an int (that
hardly anyone cared about), newer versions return a void.

\noindent{\bf Add the display of messages in the {\tt probe}, {\tt init}, {\tt exit} and
{\tt remove} methods, and observe in which order they are called, most significantly the 
{\tt probe} function when the platform\_device is created and calls the associated driver.} 

We then add pseudo-files for communicating with the user, accessible at\\
\verb!/sys/bus/platform/drivers/gpio-simple!

This result is achieved by adding a structure holding the pointer to the communication
functions and initialized with the {\tt DEVICE\_ATTR} macro.

\lstinputlisting[label=sim4,language=C]{programs/3to4}

The code becomes more challenging to read unless analyzing the source code of\\
{\tt .../buildroot/output/build/linux-headers-XXX/include/linux/device.h} and observing that
the macro is indeed defined as

\begin{verbatim}
#define DEVICE_ATTR(_name, _mode, _show, _store) \
        struct device_attribute dev_attr_##_name = __ATTR(_name, _mode, _show, _store)
\end{verbatim}

Hence, this macro creates in our case a structure named {\tt dev\_attr\_value} ({\tt value} being the
name provided as first argument to the macro) provided as argument to the {\tt 
device\_create\_file()} function.

Finally, we add the timer handler and trigger the associated event with:

\begin{lstlisting}[label=sim5,language=C]
static struct work_struct work;

static void gpio_simple_notify(struct work_struct *ws)
{pr_info("IRQ %d triggered on GPIO %d, value=%d\n",
}

static void do_something(struct timer_list *t)
{printk(KERN_INFO "plop: %lu",jiffies);
 schedule_work(&work);
}

static int gpio_simple_probe(struct platform_device *pdev)
{[...]
 INIT_WORK(&work, gpio_simple_notify);
}

static void gpio_simple_remove(struct platform_device *pdev)           
{device_remove_file(&pdev->dev, &dev_attr_value);
 cancel_work_sync(&work);
  [...]
}
\end{lstlisting}

\noindent as well as the communication with the entry point {\tt value} in the associated {\tt /sys} 
filesystem. In this example, a task is triggered when an interrupt is triggered. This event
will unlock the read function called by a userspace process when an interrupt will be
triggered.

\noindent{\bf Demonstrate how a file being read in the {\tt sysfs} is unblocked every
time the timer reaches its threshold and launches a tasklet in charge of unlocking
a semaphore.}

\subsection{Automatically loading a devicetree compatible module at boot}

We notice that when booting the Red Pitaya, if the LED or GPIO drivers were selected as
modules, they are automatically loaded after booting the board

{\footnotesize
\begin{verbatim}
redpitaya> lsmod
Module                  Size  Used by
gpio_zynq              16384  2 
leds_gpio              12288  0 
led_class              16384  1 leds_gpio
\end{verbatim}
}

Indeed, these hardware are described in the devicetree as nodes found in

{\footnotesize
\begin{verbatim}
  gpio-leds {
    compatible = "gpio-leds";
 
    led-8-yellow {
      label = "led8";
      gpios = <0x14 0x00 0x00>;
      default-state = "off";
      linux,default-trigger = "mmc0";
    };
 
    led-9-red {
      label = "led9";
      gpios = <0x14 0x07 0x00>;
      default-state = "off";
      linux,default-trigger = "heartbeat";
    };
  };
\end{verbatim}
}

\noindent so that the drivers compatible with {\tt gpio-leds} were loaded at boot time. 

The link between the {\tt compatible} keyword and the module is found in the {\tt modules.alias} of 
the {\tt /lib/modules} directory with the kernel version. Assuming we have selected to develop a 
devicetree-compatible kernel module using 

{\footnotesize
\begin{verbatim}
  static struct of_device_id simple_driver_of_match[] = {
    { .compatible = "simple_driver", }, {},
  };
MODULE_DEVICE_TABLE(of, simple_driver_of_match);
\end{verbatim}
}

\noindent then by adding the devicetree entry

{\footnotesize
\begin{verbatim}
  mynode {compatible = "simple_driver";};
\end{verbatim}
}

% Following paragraph : the {\tt modules.alias.bin} found in the {\tt /lib/modules/VER/} directory of the Red Pitaya running Linux version {\tt VER} ??? but ??? be filled
running Linux version {\tt VER} but be filled
\noindent the driver should be loaded at boot time. However, for the Linux kernel to be
aware of the link between the kernel module name and the associated {\tt compatible} keyword,
the {\tt modules.alias.bin} found in the {\tt /lib/modules/VER/} directory of the Red Pitaya
running Linux version {\tt VER} but be filled. To generated this binary file and the associated
human readable {\tt modules.alias} files, the following step are used, assuming the Buildroot
cross-compilation toolchain is locate in {\tt \$BR}:
\begin{enumerate}
\item copy the compile {\tt .ko} kernel module in {\tt \$BR/output/target/lib/modules/VER/kernel/drivers}, e.g. \\
{\tt cp flie.ko \$BR/output/target/lib/modules/6.6.0-xilinx/kernel/drivers/}
\item {\tt sudo depmod -b \$BR/output/target/ VER} so in our case\\
{\tt sudo depmod -b \$BR/output/target/ 6.6.0-xilinx}
\item check that the {\tt \$BR/output/target/lib/VER/6.6.0-xilinx/modules.alias} now ends with\\
{\tt alias of:N*T*Csimple\_driverC* file}\\
{\tt alias of:N*T*Csimple\_driver file}\\
indicating that the kernel is now aware that the {\tt compatible} keyword {\tt simple\_driver} is
associated with the {\tt file.ko} kernel module
\item copy to the Red Pitaya the content of {\tt \$BR/output/target/lib/modules/VER} or more
specifically the {\tt module.*} and the {\tt .ko} files.
\end{enumerate}

If instead of modifying the devicetree source, compiling and rebooting we wish to use a devicetree
overlay, the following overlay source

{\footnotesize
\begin{verbatim}
/dts-v1/;
/plugin/;
  
/ {
    fragment@0 {
        target-path="/";
        __overlay__ {
            demo { compatible = "simple_driver"; };
        };
    };
  
};
\end{verbatim}
}

\noindent achieves the same result of loading automatically the kernel module whose information
was provided in the {\tt module.alias} file. Notice that for the root directory, the keyword
is {\tt target-path} and not {\tt target}. We check that the new node {\tt demo} was indeed 
identified by the devicetree overlay with

{\footnotesize
\begin{verbatim}
redpitaya> cat /proc/device-tree/demo/compatible 
simple_driver
\end{verbatim}
}

\section{Debugging the kernel and kernel modules: KGDB}

The kernel debugger\footnote{\url{https://www.nxp.com/docs/en/application-note/AN14680.pdf}}\footnote{\url{https://www.linuxpromagazine.com/index.php/layout/set/print/Online/Features/Qemu-and-the-Kernel}}
must be enabled in the Buildroot Linux kernel configuration, whose menus are
reached with {\tt make linux-menuconfig} on the host buildroot top directory. KGDB can be launched
at boot time by adding {\tt rootwait} to the kernel launch command line, waiting for a GDB client 
to connect from the host as the target kernel launch is stopped, or KGDB can halt execution while 
running. The latter strategy will be used here. 

From {\tt dmesg} on the running Red Pitaya, we learn that the serial port communicating with the
console is located at {\tt ttyPS0}, so that linking KGDB with this communication interface is 
achieved with 

\noindent
\begin{verbatim}
redpitaya> echo ttyPS0 > /sys/module/kgdboc/parameters/kgdboc 
\end{verbatim}

Once the interface has been set, KGDB is launched with

\noindent
\begin{verbatim}
redpitaya> echo g > /proc/sysrq-trigger 
\end{verbatim}

The Red Pitaya halts as the kernel waits for a GDB client to connect from the host over the serial 
communication line. In order to debug the kernel, the host must hold a copy of the kernel running
on the target, as found in the Buildroot tree structure under the {\tt output/build/linux*} subdirectory:

\begin{verbatim}
$ cp $BR/output/build/linux-xilinx-v2024.1/vmlinux /tmp/
$ gdb-multiarch vmlinux
(gdb) set serial baud 115200
(gdb) target remote /dev/ttyUSB0
\end{verbatim}

\noindent assuming the serial port communicating from the host to the Red Pitaya is {\tt /dev/ttyUSB0}.
All usual commands from GDB can be used from this client. After quitting the client, the kernel execution
resumes and indeed {\tt dmesg} on the Red Pitaya confirms that {\tt KGDB: Entering KGDB} the debugger
had been launched.

This allows for debugging the Linux kernel, but debugging a kernel module is slightly more involved since
loading the init function must be intercepted by setting a breakpoint at the appropriate location of 
the kernel. The location of the {\tt find\_module\_section} function is located by searching the Linux
kernel source code on the host Buildroot:

\noindent
\begin{verbatim}
grep -A2 -n find_module_section $BR/output/build/linux-xilinx-v2024.1/kernel/module/main.c 
2903:   err = find_module_sections(mod, info);
2904-   if (err)
2905-           goto free_unload;
\end{verbatim}

\noindent and the line following the {\tt err = find\_module\_section} will define the breakpoint location. Setting the
breakpoint is achieved with

\begin{verbatim}
$ gdb-multiarch vmlinux 
(gdb) set serial baud 115200
(gdb) target remote /dev/ttyUSB0
(gdb) break kernel/module/main.c:2903
Breakpoint 1 at 0xc0175f5c: file kernel/module/main.c, line 2903.
(gdb) continue
Continuing.
\end{verbatim}

We wish to intercept the init function of a kernel module named {\tt 0mymod.ko}. Doing so requires knowing some
memory organization:

\begin{verbatim}
redpitaya> insmod /tmp/0mymod.ko 
redpitaya> cat /sys/module/0mymod/sections/.init.text 
0xbf014000
\end{verbatim}

Once we know where the init function memory location is, we remove the module from memory ({\tt rmmod}) and resume
KGDB:

\begin{verbatim}
redpitaya> echo g > /proc/sysrq-trigger 
\end{verbatim}

\noindent and connect a GDB client to run the command {\tt continue} and resum kernel execution. Upon loading
the kernel module, this time

\begin{verbatim}
redpitaya> insmod /tmp/0mymod.ko 
[New Thread 190]
[Switching to Thread 192]
 
Thread 69 hit Breakpoint 1, find_module_sections (info=0xe0929ef0, mod=<optimized out>) at
	kernel/module/main.c:2104
2104            mod->kp = section_objs(info, "__param",
\end{verbatim}

\noindent and on the host:

\begin{verbatim}
(gdb) add-symbol-file 0mymod.ko 0xbf014000
(gdb) continue
Thread 67 hit Breakpoint 2.2, hello_end () at .../0mymod.c:6
6       static void __exit hello_end(void) {printk(KERN_INFO "Goodbye\n");}
\end{verbatim}

\noindent giving access to the kernel module memory location, including the init function.

%(gdb) print mod->init_layout.base
%value has been optimized out

%Make sure to enable (/kallsy)
%[*] General setup -> Load all symbols for debugging / ksymoops -> Include all symbols in kallsyms
%Symbol: KALLSYMS_ALL [=y]

% make V=1

% debug symbols:
% https://stackoverflow.com/questions/28298220/kernel-module-no-debugging-symbols-found

% Twice Linux ???
\section{Compiling a Linux image Linux with Xenomai support (real-time extension but no DTBO 
support)}\label{xe}

Xenomai has been ported to the Zynq
% \footnote{\url{https://session.wikispaces.net/57898/auth/auth?authToken=0eeca7316e2df30bb5b70760aa5847fe8} 
and most significantly J.H. Brown \& B. Martin, {\em How fast is fast enough? Choosing between 
Xenomai and Linux for real-time applications}, Proc. 12th Real-Time Linux 
Workshop (RTLWS’12) (2012)}. The updates brought to an official mainline kernel were adapted
to the {\tt linux-xlnx} kernel used here and included as patches to be applied when compiling
Xenomai from Buildroot.

%\noindent\fbox{\parbox{\linewidth}{Default patches included in Buildroot for compiling
%Xenomai 3.0.5 {\bf must be deleted}: these patches conflict with our own updates to the
%{\tt linux-xlnx} kernel. Hence, in the buildroot directory, {\tt rm packages/xenomai/*.patch}
%will remove the unnecessary patches.}}
%
%The reader willing to modify an official mainline kernel can follow the procedure described at
%\url{https://embeddedgreg.com/2017/07/16/getting-started-with-xenomai-3-on-zynq/comment-page-1/},
%considering that using such a mainline kernel will prevent support for accessing the
%FPGA from GNU/Linux. This procedure should {\bf not} be followed according to the steps
%described in the usage of buildroot as follows.

Configuring buildroot with Xenomai support for the Red Pitaya is called with the
{\tt redpitaya\_xenomai\_defconfig} configuration file. Hence, in the buildroot directory
and after exporting the {\tt BR2\_EXTERNAL} variable towards the {\tt redpitaya} support files,
run the {\tt make redpitaya\_xenomai\_defconfig} command to load the configuration, which is
then applied during the compilation with {\tt make}. As before, after quite some time and
additional storage space use, the resulting image is located in {\tt output/images} and is
transferred to the SD card with the appropriate {\tt dd} command.

In case the FPGA (PL) of the Zynq is to be used, make sure to activate its support in the kernel
with {\tt make linux-menuconfig} and after searching for FPGA, select {\em FPGA Configuration 
Framework} and in the sub-menu:

\begin{verbatim}
< >   FPGA debug fs
< >   Altera Partial Reconfiguration IP Core
< >   Altera FPGA Passive Serial over SPI 
< >   Altera CvP FPGA Manager 
<*>   Xilinx Zynq FPGA      
< >   Xilinx Configuration over Slave Serial (SPI)
< >   Lattice iCE40 SPI
< >   Lattice MachXO2 SPI
<*>   FPGA Bridge Framework
< >     Altera FPGA Freeze Bridge
< >     Altera FPGA Freeze Bridge
<*>   Xilinx LogiCORE PR Decoupler
<*>   FPGA Region
<*>     FPGA Region Device Tree Overlay Support
< >   FPGA Device Feature List (DFL) support
\end{verbatim}

as well as activating

\begin{verbatim}
Device Tree and Open Firmware support
-*-   Device Tree overlays
\end{verbatim}

%%En particulier, nous devons
%\begin{enumerate}
%\item t\'el\'echarger le noyau de la branche officielle par \\
%{\tt git clone git://git.kernel.org/pub/scm/linux/kernel/git/stable/linux-stable.git} 
%et passer dans la version 4.9.24 par \\
%{\tt git checkout tags/v4.9.24 -b zynq\_xeno\_4.9.24}
%\item t\'el\'echarger Xenomai par {\tt git clone git://git.xenomai.org/xenomai-3.git} et
%passer dans la branche 3.0.5 par {\tt git checkout tags/v3.0.5 -b xeno\_3.0.5} dans le
%r\'epertoire de travail ainsi cr\'e\'e
%\item t\'el\'echarger le patch ipipe au m\^eme niveau d'arborescence que le noyau Linux et 
%Xenomai~: \\
%{\tt wget https://xenomai.org/downloads/ipipe/v4.x/arm/ipipe-core-4.9.24-arm-2.patch}
%\item compiler Xenomai en entrant dans son r\'epertoire et en faisant appel au patch que nous
%venons de t\'el\'echarger~:\\
%{\tt ./scripts/prepare-kernel.sh --arch=arm --ipipe=../ipipe-core-4.9.24-arm-2.patch --linux=\$LINUX\_HOME}
%avec la variable pointant vers le r\'epertoire contenant les sources du noyau t\'el\'echarg\'ees
%au cours de la premi\`ere \'etape
%\item maintenant que le noyau est modifi\'e pour supporter son insertion dans Xenomai, il faut
%le configurer et le compiler. Entrer dans le r\'epertoire du noyau ({\tt cd \$LINUX\_HOME}) et
%copier la configuration pour processeur Zynq, disponible par exemple \`a
%\url{https://raw.githubusercontent.com/Xilinx/linux-xlnx/master/arch/arm/configs/xilinx\_zynq\_defconfig
%}, dans le r\'epertoire {\tt arch/arm/configs/xilinx\_zynq\_defconfig} des sources de Linux. 
%Prendre en compte cette configuration par\\
%{\tt make ARCH=arm xilinx\_zynq\_defconfig} \\
%puis corriger quelques d\'eficiences du noyau pour tenir compte de l'utilisation dans un contexte
%d'application temps-r\'eelle en retirant la capacit\'e \`a modifier la vitesse du processeur ainsi
%que l'organisation de la m\'emoire~:\\
%{\tt make ARCH=arm menuconfig}\\
%d\'esactiver {\tt CPU Power Management} $\rightarrow$ {\tt CPU Frequency Scaling}, ainsi que
%{\tt Kernel Features} $\rightarrow$ {\tt contiguous memory allocator}. Nous n'avons pas rencontr\'e
%le probl\`eme mentionn\'e concernant l'option {\tt irq tracing}.
%\item nous v\'erifions que ce nouveau noyau compile bien par \\
%{\tt make ARCH=arm CROSS\_COMPILE=arm-buildroot-linux-uclibcgnueabihf- UIMAGE\_LOADADDR=0x8000 uImage modules dtbs} \\
%en ayant pris soin d'ins\'erer le compilateur fourni par {\tt buildroot} (disponible dans
%{\tt output/host/usr/bin} de {\tt buildroot}) dans le chemin d'ex\'ecutables ({\tt PATH}).
%\item Si tout se passe bien et qu'un image uImage est g\'en\'er\'ee, nous pouvons ins\'erer
%ce noyau dans le flux de traitement {\tt buildroot}. Pour ce faire, nous nettoyons les sources du
%noyau des binaires g\'en\'er\'es lors de la compilation ({\tt make clean}), nous fabriquons une
%archive des sources du noyau auquel nous avons appliqu\'e les modifications Xenomai ({\tt
%tar zcvf /tmp/kernel.tar.gz linux-stable/*}) et nous expliquons \`a {\tt buildroot}
%d'utiliser ce noyau pour g\'en\'erer l'image qui sera flash\'ee sur carte SD pour \^etre
%ex\'ecut\'ee sur Zynq~: dans {\tt buildroot}, {\tt Kernel} $\rightarrow$ {\tt Kernel version} 
%$\rightarrow$ {\tt Custom tarball} et {\tt Kernel} $\rightarrow$ {\tt URL of custom kernel tarball}
%$\rightarrow$ \url{file:///tmp/kernel.tar.gz} et retirer le contenu de {\tt Kernel} $\rightarrow$ 
%{\tt Custom kernel patches}
%\item suite \`a ces modifications, {\tt make -j8} dans {\tt buildroot} recompile une image
%contenant un noyau Linux avec les fonctions temps r\'eel de Xenomai.
%\end{enumerate}

Following the compilation and execution of this new image, new messages are displayed upon
booting Linux confirming that Xenomai and associated tools have been included in the kernel:
\begin{verbatim}
sched_clock: 64 bits at 333MHz, resolution 3ns, wraps every 4398046511103ns
I-pipe, 333.333 MHz clocksource, wrap in 12884 ms
\end{verbatim}
and
\begin{verbatim}
hw perfevents: enabled with armv7_cortex_a9 PMU driver, 7 counters available
[Xenomai] scheduling class idle registered.
[Xenomai] scheduling class rt registered.
I-pipe: head domain Xenomai registered.
[Xenomai] Cobalt v3.0.5 (Sisyphus's Boulder) 
workingset: timestamp_bits=30 max_order=17 bucket_order=0
\end{verbatim}

%Pour obtenir les programmes utilisateurs pour stresser le syst\`eme et qualifier le bon
%fonctionnement de Xenomai, entrer dans le r\'epertoire Xenomai t\'el\'echarg\'e
%ci-dessus ({\tt cd [...]/xenomai-3}) et configurer l'environnement de travail par\\
%{\tt ./scripts/bootstrap}. Comme souvent maintenant, il est peu judicieux de compiler
%directement dans le r\'epertoire des sources~: nous cr\'eons donc un r\'epertoire pour
%compiler les ex\'ecutables ({\tt mkdir build}) dans lequel nous entrons ({\tt cd build})
%avant de configurer les outils de Xenomai pour \^etre compil\'es avec la cha\^\i ne de compilation
%fournie par {\tt buildroot}~:
%\begin{verbatim}
%../configure CFLAGS="-march=armv7-a -mfpu=vfp3 -mfloat-abi=hard" LDFLAGS="-march=armv7-a" \
%--build=i686-pc-linux-gnu --host=arm-none-linux-gnueabi --with-core=cobalt \
%--enable-smp --enable-tls CC=arm-buildroot-linux-uclibcgnueabihf-gcc \
%LD=arm-buildroot-linux-uclibcgnueabihf-ld
%\end{verbatim}
%et finalement la compilation s'ach\`eve par \\
%\verb~make DESTDIR=/tmp/xeno_zynq/ install~
%qui cr\'ee une arborescence {\tt xeno\_zynq} dans le r\'epertoire temporaire {\tt /tmp}. Tous
%les binaires dans ce r\'epertoire sont \'evidemment \`a destination du processeur ARM, donc nous
%copions l'int\'egralit\'e de son contenu vers la Red Pitaya ({\tt scp -r /tmp/xeno\_zynq
%root@192.168.2.200:/tmp})
%
%Nous nous connectons sur la Red Pitaya pour v\'erifier que nos outils fonctionnent. Pour ce faire,
%nous exportons le chemin vers les biblioth\`eques dynamiques
%\begin{verbatim}
%redpitaya> export LD_LIBRARY_PATH=/tmp/xeno_zynq/usr/xenomai/lib/
%\end{verbatim}
%et 

A classical tool is provided by Xenomai called {\tt latency} for qualifying delays of the 
real-time operating system, in which the generic Linux kernel has only become one additional 
non-real time task:
% ex\'ecutons un des exemples fourni dans {\tt /tmp/xeno\_zynq/usr/xenomai/demo} ou 
\begin{verbatim}
redpitaya> latency 
== Sampling period: 1000 us
== Test mode: periodic user-mode task
== All results in microseconds
warming up...
RTT|  00:00:01  (periodic user-mode task, 1000 us period, priority 99)
RTH|----lat min|----lat avg|----lat max|-overrun|---msw|---lat best|--lat worst
RTD|     -2.798|     -2.435|      2.714|       0|     0|     -2.798|      2.714
RTD|     -2.811|     -2.374|      7.931|       0|     0|     -2.811|      7.931
RTD|     -2.380|     -1.888|      6.692|       0|     0|     -2.811|      7.931
RTD|     -2.806|     -2.393|      6.939|       0|     0|     -2.811|      7.931
RTD|     -2.836|     -2.397|      5.680|       0|     0|     -2.836|      7.931
RTD|     -2.870|     -2.358|      7.046|       0|     0|     -2.870|      7.931
RTD|     -2.883|     -2.388|      8.137|       0|     0|     -2.883|      8.137
\end{verbatim}

Negative values -- meaningless -- are documented at \\
\url{https://embeddedgreg.com/2017/07/16/getting-started-with-xenomai-3-on-zynq/comment-page-1/} 
and we indeed find on the Red Pitaya the pseudo-filesystem {\tt /proc/xenomai}, as well as the
{\tt autotune} tool in \\{\tt /tmp/xeno\_zynq/usr/xenomai/sbin} used for tuning the system
parameters \footnote{{\tt autotune --period 10000}}, in order to finally achieve a meaningful
result:
\begin{verbatim}
redpitaya> latency 
== Sampling period: 1000 us
== Test mode: periodic user-mode task
== All results in microseconds
warming up...
RTT|  00:00:01  (periodic user-mode task, 1000 us period, priority 99)
RTH|----lat min|----lat avg|----lat max|-overrun|---msw|---lat best|--lat worst
RTD|      0.993|     10.388|     26.247|       0|     0|      0.993|     26.247
RTD|      0.980|      8.411|     29.984|       0|     0|      0.980|     29.984
RTD|      1.137|      9.887|     25.449|       0|     0|      0.980|     29.984
RTD|      1.180|     11.129|     26.790|       0|     0|      0.980|     29.984
RTD|      1.152|     10.576|     28.543|       0|     0|      0.980|     29.984
RTD|      1.055|     10.865|     27.385|       0|     0|      0.980|     29.984
RTD|      1.063|     10.913|     28.225|       0|     0|      0.980|     29.984
\end{verbatim}

Another tool called {\tt stress} allows to load the system and observe that latencies remain
stable despite the resources requested by this program.
% Comme quelques outils de Xenomai s'attendent \`a trouver les binaires dans {\tt /usr/xenomai},
%nous pouvons les satisfaire par un lien symbolique 
%\begin{verbatim}
%redpitaya> ln -s /tmp/xeno_zynq/usr/xenomai /usr/xenomai
%\end{verbatim}
%afin de rendre par exemple {\tt xeno-test-run-wrapper} heureux.

%{\bf JMF : A FINIR -- AU LIEU DE COMPILER LE KERNEL XENO EN EXTERNE, INSERER DANS LE FLUX BR, ET
%COMMENT RETROUVER XDEVCFG DANS UN KERNEL MAINLINE + REFAIRE LES EXEMPLES GPIO
%
%pourquoi testsuite de Xenomai dans buildroot ne compile pas ? buildroot essaie d'appliquer
%une 2eme fois un patch qui n'a pas de sens

\begin{itemize}
\item check on your system that the real time options have been activated, and most significantly
the associated tools ({\em testsuite} and {\tt rt-tests}). Achieving this result requires
activating {\tt Target packages} $\rightarrow$ {\tt Real-Time} $\rightarrow$ {\tt Xenomai Userspace} 
$\rightarrow$ {\tt Install testsuite}. Furthermore, manually select the Xenomai version by filling
the {\tt Custom Xenomai version} field with {\tt 3.0.5},
\item some additional shell tools might be wisely added, for example {\tt screen} in {\tt Target 
packages} $\rightarrow$ {\tt Shell and utilities} $\rightarrow$ {\tt screen} or {\tt Target 
packages} $\rightarrow$ {\tt System tools} $\rightarrow$ {\tt cpuload}. 
\end{itemize}

The latencies of the real time system \cite{blaes} are qualified with
\colorbox{lightgray}{\verb~echo "0" > /proc/xenomai/latency~} followed with
\colorbox{lightgray}{\verb~latency -p 100~}. Observe the latency fluctuations when, in another
terminal, the processor is loaded with\\
\colorbox{lightgray}{\verb~while ( true ) ; do echo toto;done~}. This last command is run in a second
shell started from {\tt screen} (CTRL-A c to create a new shell session, CTRL-A a to switch between
sessions).

%{\footnotesize
%\begin{verbatim}
%# latency -p 100
%== Sampling period: 100 us
%== Test mode: periodic user-mode task
%== All results in microseconds
%warming up...
%RTT|  00:00:01  (periodic user-mode task, 100 us period, priority 99)
%RTH|----lat min|----lat avg|----lat max|-overrun|---msw|---lat best|--lat worst
%RTD|      5.083|      5.166|     11.583|       0|     0|      5.083|     11.583
%RTD|      5.749|      5.874|     14.124|       0|     0|      5.083|     14.124
%RTD|      5.708|      5.874|     11.083|       0|     0|      5.083|     14.124
%RTD|      5.708|      5.874|     10.999|       0|     0|      5.083|     14.124
%RTD|      5.749|     10.374|     18.999|       0|     0|      5.083|     18.999  <- lancement dans un
%RTD|      6.416|     10.624|     18.624|       0|     0|      5.083|     18.999  <- second terminal de
%RTD|      6.291|     10.624|     18.291|       0|     0|      5.083|     18.999  <-  while ( true ) ; do
%RTD|      6.124|     10.666|     18.249|       0|     0|      5.083|     18.999  <-    echo toto;done
%RTD|      5.666|      7.499|     17.458|       0|     0|      5.083|     18.999 
%RTD|      5.708|      5.874|     12.124|       0|     0|      5.083|     18.999
%RTD|      5.749|      5.874|     11.291|       0|     0|      5.083|     18.999
%\end{verbatim}
%}
%
%\section{Cross compiler~: clignotement d'une LED}
%
%Cross-compiler un programme pour Olinuxino-A13 depuis un PC revient \`a utiliser le 
%compilateur appropri\'e pour g\'en\'erer un binaire compr\'ehensible par un processeur 
%d'architecture ARM. Afin de faire clignoter une diode, nous sommes inform\'es \`a
%\url{http://olimex.wordpress.com/2012/10/23/a13-olinuxino-playing-with-gpios/} que
%la diode verte de statut se trouve sur le port GPIO\_G9. Une biblioth\`eque
%en espace utilisateur acc\'edant directement aux adresses des ports (sans passer par
%le noyau Linux) est mise \`a disposition \`a \url{http://docs.cubieboard.org/tutorials/common/gpio_on_lubuntu}~:
%nous cross-compilons un exemple trivial d'allumage et d'extinction de diode nomm\'e {\tt example.c}
%
%\lstinputlisting[language=C]{../gpio/gpio_sleep.c}
%
%par
%
%\begin{verbatim}
%arm-linux-gnueabihf-gcc -c example.c
%arm-linux-gnueabihf-gcc -c gpio_lib.c
%arm-linux-gnueabihf-gcc -o example example.o gpio_lib.o 
%\end{verbatim}
%
%{\bf Envoyer le fichier binaire depuis le PC vers la carte Olinuxino et constater le r\'esultat. Reproduire
%en encodant par {\tt uuencode}}.
%
%En cas d'absence de certaines biblioth\`eques ({\tt Command not found}), ajouter dans le r\'epertoire
%{\tt lib} de la seconde partition de la carte SD contenant le syst\`eme d'exploitation les fichiers {\tt
%libc.so.6} et {\tt ld-linux.so.3} (cf le r\'esultat de {\tt ldd example}).
%
%Ayant valid\'e la capacit\'e \`a commander l'allumage et l'extinction d'une diode depuis l'espace
%utilisateur, nous allons exploiter cette fonctionnalit\'e pour comparer les latences lors de
%l'ex\'ecution d'une t\^ache p\'eriodique avec ou sans l'option temps-r\'eel. Les quatre options qui
%s'offrent \`a nous, excluant la solution d'un module noyau, sont des temporisations par {\tt sleep} ou
%par {\em timer}, avec ou sans l'extension temps-r\'eel de Xenomai. Dans tous les cas, nous
%lan\c cons sur l'Olinuxino-A13-micro, auquel nous nous connectons par le port s\'erie virtuel au moyen
%de {\tt kermit -c}, une session {\tt screen}. Un terminal sert \`a lancer le programme commutant 
%l'\'etat de la diode de statut (verte) tandis que le second terminal charge ou non le processeur
%en ex\'ecutant {\tt while (true); do echo toto;done}.

\section{Latency qualification}

% Next paragraph : Does ``explicited'' exist in English ???
The four examples explicited below investigate how different delays between toggling
LED state either by calling the {\tt sleep} function (active delay), or by calling a timer
callback function. Without Xenomai, calling signal callbacks leads to the worst stability
as soon as the processor becomes loaded. In all other cases, quantifying latency is achieved
by triggering an oscilloscope on the rising edge of the LED state and watching the falling edge
delay. If the delay met exactly the expected duration, all falling edges would overlap. The
time interval over which the falling edges are observed quantifies the resistance of each method
to processing load and the ability to meet the expected latency requirements under the various
delay conditions. The case of a kernel module is not addressed.

% \newpage

\subsection{Delay using {\tt sleep}, without Xenomai}

\begin{minipage}[t]{\linewidth}
\begin{minipage}{0.59\linewidth}
\lstinputlisting[language=C]{./programs/gpio_sleep.c}
\end{minipage}
\begin{minipage}{0.39\linewidth}
\includegraphics[width=\linewidth]{sleep_periode.png} \\
Period\\
\includegraphics[width=\linewidth]{sleep_charge.png} \\
Loaded\\
\includegraphics[width=\linewidth]{sleep_vide.png}\\
Unloaded
\end{minipage}
\end{minipage}

This program follows the previous example seen earlier of toggling LED state, with a 500~ms delay.
In this and all the following examples, the processor is loaded with the following command executed in
a second terminal on the Red Pitaya:
\verb~stress --cpu 8 --io 4 --vm 2 --vm-bytes 128M --timeout 10s~ executed 5 times in a row for a total
measurement duration of about a minute.

\subsection{Delay using a timer, without Xenomai}

\begin{minipage}[t]{\linewidth}
\begin{minipage}{0.59\linewidth}
\lstinputlisting[language=C]{./programs/gpio_sigalarm.c}
\end{minipage}
\begin{minipage}{0.39\linewidth}
\includegraphics[width=\linewidth]{sigalarm_periode.png} \\
Period\\
\includegraphics[width=\linewidth]{sigalarm_charge.png} \\
Loaded\\
\includegraphics[width=\linewidth]{sigalarm_vide.png}\\
Unloaded
\end{minipage}
\end{minipage}

The result is catastrophic: no only is sigalarm unable to meet latency constraints, but
even when unloaded periodic delay jumps are observed. This delay method is unable to cope
with processor load.

\subsection{Delay using {\tt sleep}, with Xenomai}

\begin{minipage}[t]{\linewidth}
\begin{minipage}{0.59\linewidth}
\lstinputlisting[language=C]{./programs/xeno_gpio_sleep.c}
\end{minipage}
\begin{minipage}{0.39\linewidth}
\includegraphics[width=\linewidth]{xeno_sleep_periode.png} \\
Period\\
\includegraphics[width=\linewidth]{xeno_sleep_charge.png} \\
Loaded\\
\includegraphics[width=\linewidth]{xeno_sleep_vide.png}\\
Unloaded
\end{minipage}
\end{minipage}

\subsection{Delay using a timer, with Xenomai}

\begin{minipage}[t]{\linewidth}
\begin{minipage}{0.59\linewidth}
\lstinputlisting[language=C]{./programs/xeno_gpio_timer.c}
\end{minipage}
\begin{minipage}{0.39\linewidth}
\includegraphics[width=\linewidth]{xeno_timer_periode.png} \\
Period\\
\includegraphics[width=\linewidth]{xeno_timer_charge.png} \\
Loaded\\
\includegraphics[width=\linewidth]{xeno_timer_vide.png}\\
Unloaded
\end{minipage}
\end{minipage}

%sigalarm_charge.png
%sigalarm_periode.png
%sigalarm_vide.png
%sleep_charge.png
%sleep_periode.png
%sleep_vide.png
%xeno_sleep_charge.png
%xeno_sleep_periode.png
%xeno_sleep_vide.png

\section{FPGA configuration}

% Next paragraph : accessined ???
While a classical FPGA requires a dedicated framework for configuration (e.g. JTAG probe),
the Zynq allows accessined the Programmable Logic part (PL) of the chip from its general
purpose Processing System (PS). 
%Two methods are provided to achieve such a result:
%historicaly, Xilinx has added a driver module creating {\tt /dev/xdevcfg} in which writing
%({\tt cat}) the bitstram to this communication interface would configure the FPGA. 
%However, in an attempt to provide an homogeneous interface for all System On Chips (SoC) --
%including Altera/Intel and Lattice/Microsemi -- a new 
A framework called FPGA Manager has been 
standardized. The benefit of this framework is that it applies to {\em devicetree overlays}
aimed a synchronizing all resources needed by a bitstream used for configuring the FPGA
as well as the associated Linux driver. Using the devicetree overlay, the kernel module
with the driver is automatically loaded when the associated bitstream defined in the
plugin is requested, creating the link between the FPGA, the kernel module running on the processing
system, and the userspace Linux.

%With the proposed kernel, both approaches can be used. Due to conflicting access to the unique
%PL resource, only one method can be active at any given time. The order in which kernel modules
%are loaded during the boot sequence leads the FPGA Manager to be the default option activated
%when starting Linux. Activating {\tt /dev/xdevcfg}, requires removing all modules related
%to FPGA Manager and inserting the {\tt xilinx\_devcfg} module which creates {\tt /dev/xdevcfg}: 
%when booting, the system logs ({\tt dmesg}) the activation of FPGA Manager with the message:
%\begin{verbatim}
%FPGA manager framework
%fpga_manager fpga0: Xilinx Zynq FPGA Manager registered
%fpga-region fpga-full: FPGA Region probed
%random: crng init done
%\end{verbatim}
%
%All associated modules are removed and the {\tt xilinx\_devcfg} module is inserted:
%\begin{verbatim}
%rmmod zynq_fpga
%rmmod xilinx_devcfg
%modprobe xilinx_devcfg
%\end{verbatim}
%
%leading to 
%\begin{verbatim}
%fpga_manager fpga0: fpga_mgr_unregister Xilinx Zynq FPGA Manager
%xdevcfg f8007000.devcfg: ioremap 0xf8007000 to dea36000
%\end{verbatim}
%
%and indeed {\tt /dev/xdevcfg} has been created:
%\begin{verbatim}
%redpitaya> ls -l /dev/xdevcfg 
%crw-rw----    1 root     root      240,   0 Jan  1 13:47 /dev/xdevcfg
%\end{verbatim}
%
%Nous revenons \`a FPGA Manager par la proc\'edure inverse~:
%\begin{verbatim}
%redpitaya> rmmod xilinx_devcfg
%redpitaya> modprobe zynq_fpga
%redpitaya> ls -l /sys/class/fpga_manager/fpga0/
%[...]
%--w-------    1 root     root          4096 Jan  1 13:52 firmware
%\end{verbatim}

\input{fpga_manager_zynq_eng.tex}

\section{PS-PL interaction}

\subsection{GPIO}

The major benefit of the Zynq architecture is to combine a general purpose Processing
System (PS) processor executing GNU/Linux, and a reconfigurable Programmable Logic (PL) 
FPGA. We will demonstrate in the following example how to synthesize the GPIO IP 
configuring the PL, unlike the MIO we used earlier that can be accessed only by the PS, 
and how to access this new peripheral using the techniques developed in this document:
\begin{itemize}
\item direct access from userspace
\item direct access from kernel space (module)
\item accessing from a kernel driver probed by the {\em devicetree} using the {\tt compatible} method.
\end{itemize}

Configuring the PL from Vivado is beyond the scope of this document. We here use a GPIO block provided
by Xilinx as a Vivado IP, and connect to the PS through the AXI bus, a communication standard used on all
ARM-processor based platforms (Fig. \ref{vivado}).

\begin{figure}[h!tb]
\includegraphics[width=0.65\linewidth]{vivado.png}
\includegraphics[width=0.34\linewidth]{axi_gpio.png}
\caption{The PL configured with an AXI Interconnect to link the GPIO IP in the FPGA with the PS.}
\label{vivado}
\end{figure}

As any peripheral on digital systems, the GPIO is connected through the three interfaces of the
AXI bus, namely the data bus, control bus and address bus. The logic of handling the control
bus, quite complex, is hidden with the AXI Interconnect and the only information we will be interested
in is the base address allocated to the IP and the address offset of each control register. The
base address, varying with the kind and number of peripherals implemented in the FPGA, is provided
by Vivado (Fig. \ref{addr1}). The address offsets to the various register holding the configuration
functions are described in the GPIO bloc documentation provided by Xilinx
\footnote{p.10 of \url{https://docs.amd.com/r/en-US/pg144-axi-gpio/}}. 
Hence, we know how to access each PL peripheral function from the PS as we would have done on a microcontroller (Fig. 
\ref{addr1}).

\begin{figure}[h!tb]
\includegraphics[width=0.69\linewidth]{adresse.png}
\includegraphics[width=0.30\linewidth]{axi_gpio_regs.png}
\caption{Address assignment of the peripherals added to the PL.}
\label{addr1}
\end{figure}

\subsection{XADC}

The Zynq is fitted with a slow analog to digital converter. This peripheral can be accessed
either from the PS, or from the PL. Xilinx provides an IIO compatible driver as expected for
such a peripheral. Without dedicated configuration of the PL, only the temperature can be read
this way. We here wish to extend such PL functionalities by updating the PL configuration to
configure two pins for analog voltage measurements.

\begin{figure}[h!tb]
\includegraphics[width=0.53\linewidth]{config_xadc0.png}
\includegraphics[width=0.46\linewidth]{config_xadc2.png}
\caption{Steps for configuring the XADC and connect analog inputs. Connect the missing pins and run the
AXI bus autorouting after having added the XADC processing block.}
\label{xadc1}
\end{figure}

\begin{figure}[h!tb]
\begin{center}
\includegraphics[width=0.65\linewidth]{config_xadc1.png}
\end{center}
\caption{Steps for configuring the XADC and connect analog inputs. A bug in Vivado requires activating the
Channel Sequencer mode before setting the Independent ADC mode when configuring the list of active channels.}
\label{xadc2}
\end{figure}

The devicetree overlay for communicating with this new configuration of the XADC is
{\footnotesize
\begin{verbatim}
/dts-v1/;
/plugin/;
/ {compatible = "xlnx,zynq-7000";

   fragment@0 {
        target = <&fpga_full>;
        #address-cells = <1>;
	#size-cells = <1>;
	__overlay__ {
		#address-cells = <1>;
		#size-cells = <1>;
		firmware-name = "system_wrapper.bit.bin";
		};
	};
	fragment@1 {
		target = <&adc>;
		__overlay__ {
			xlnx,channels {
				#address-cells = <1>;
				#size-cells = <0>;
				channel@0 {reg = <0>;};
				channel@1 {reg = <1>;};
				channel@2 {reg = <2>;};
				channel@9 {reg = <9>;};
				channel@10 {reg = <10>;};
			};
		};
	};
};

\end{verbatim}
}

which overloads the {\tt \&adc} entry already present in the Red Pitaya devicetree. The organization
and options of the Xilinx driver for the ADC are descibed in the kernel documentation at \\
{\tt Documentation/devcietree/bindings/iio/adc/xilinx-xadc.txt}. We indeed find in \\
{\tt arch/arm/boot/dts/xilinx/zynq-7000.dtsi} of the kernel sources the configuration of a Zynq-7000 with
its ADC entry:
\begin{verbatim}
adc: adc@f8007100 {
   compatible = "xlnx,zynq-xadc-1.00.a";
   reg = <0xf8007100 0x20>;
   interrupts = <0 7 4>;
   interrupt-parent = <&intc>;
   clocks = <&clkc 12>;
};
\end{verbatim}
which is overloaded with the new configuration. The Xilinx driver, defined as compatible with 
{\tt xlnx,zynq-xadc-1.00.a}, is configured to comply with the IIO inputs and outputs as described in the documentation 
cited previously. We hence just need to modify the devicetree, without having to update the driver code to access
the new resources implemented in the bitstream including the XADC connected to the PS through the AXI bus. We
check the location of the driver module with
\begin{verbatim}
$ LINUX/drivers/iio/adc$ grep xlnx * | grep adc
xilinx-xadc-core.c:     { .compatible = "xlnx,zynq-xadc-1.00.a", (void *)&xadc_zynq_ops },
\end{verbatim}
which is indeed stored with other IIO drivers.

The base address of the XADC connected to the PS is documented at page 1160 of \url{ug585-Zynq-7000.pdf}: the registers for configuring the XADC from the PS are located at addresses starting at 0xf8007100. In order to comply with the XADC address organization in the bitstream and describe this organization in the devicetree when loading the IIO driver, we must compile the {\tt xilinx\_xadc.ko} driver as a module (as opposed to a static link into the kernel). 
% recompiler le driver XADC en MODULE afin de le recharger avec la nouvelle config fournie dans le dts
% make linux && make
Once loaded, the IIO peripheral is accessible at {\tt /sys/bus/iio/devices/iio:device0}.

The \url{Red_Pitaya_Schematics_STEM_125-10_V1.0.pdf} schematic describes the relation between the information screen printed on the printed circuit board and the Zynq pin names and their functions: AI0 is connected to AD8, AI1 is connected to AD0, AI2 is connected to AD1 and AI3 is connected to AD9, which helps understanding the naming chosen in the IIO driver:

\begin{minipage}[t]{\linewidth}
\begin{minipage}{.49\linewidth}
\begin{itemize}
\item pin 16 E2=in\_voltage8\_vaux9\_raw
\item pin 13 E2=in\_voltage9\_vaux8\_raw
\item pin 14 E2=in\_voltage11\_vaux0\_raw
\item pin 15 E2=in\_voltage10\_vaux1\_raw
\end{itemize}

In order to access the XADC registers from the PL, their address is described in the documentation
of the XADC Wizard v3.0 (reference PG091, version April 1, 2015): for example the register for accessing
ADC8 is located at offset 0x260 with respect to the base address provided by Vivado.

\end{minipage}
\begin{minipage}{.49\linewidth}
\includegraphics[width=\linewidth]{xadc_pinout.png}
\end{minipage}
\end{minipage}

\appendix
%\hrule
%\section{Portage de Xenomai, sur OlinuxIno A13 micro, dans buildroot}
%
%La mise en \oe uvre de Xenomai sur un processeur particulier passe par la
%pr\'esence, pour cette architecture, du support de \textbf{Ipipe}.
%En effet, pour autoriser
%l'ex\'ecution en parall\`ele de deux noyaux (l'un temps-r\'eel, l'autre
%temps-partag\'es), avec une garantie de latence minimale, il est n\'ecessaire
%d'ajouter une couche d'abstraction entre le mat\'eriel et le logiciel. Celle-ci
%va mettre \`a disposition, dans l'ordre des priorit\'es, les interruptions
%mat\'erielles, et \'eviter qu'un des domaines ne se trouve bloqu\'e sur
%l'attente d'une interruption masqu\'ee par l'autre domaine.
%
%Le probl\`eme qui se pose avec les processeurs allwinner est qu'il n'existe pas,
%\`a l'heure actuelle, de support officiel.
%
%Toutefois, Paul Crubille
%\footnote{\url{http://www.xenomai.org/pipermail/xenomai/2013-February/027801.html}} a r\'ealis\'e ce travail sur noyau sunxi 3.4.24 
%et le propose au travers de deux patchs~:
%\begin{enumerate}
%\item \textbf{pre.patch} qui supprime
%les modifications apport\'ees entre la version 3.4.6 (version sur lequel Ipipe
%s'applique) et la version 3.4.24 du noyau linux (en tenant compte \'egalement
%des modifications entre un noyau officiel et celui propos\'e par sunxi;
%\item \textbf{post.patch} qui re-ajoute
%les modifications supprim\'ees par \textbf{pre.patch} ainsi que le support Ipipe
%pour a13. 
%\end{enumerate}
%Il est ainsi possible, de mani\`ere manuelle, d'appliquer
%\textbf{pre.patch}, de lancer \textbf{prepare-kernel} fournit dans l'archive de
%Xenomai en lui sp\'ecifiant le patch Ipipe adapt\'e et finalement d'appliquer 
%\textbf{post.patch} afin d'obtenir un noyau avec le support temps-r\'eel pour
%les processeurs A13.
%
%\subsection{Cr\'eation d'un patch Ipipe avec le support Red Pitaya}\label{ipipe}
%
%La solution en trois patchs (pre, prepare et post), dans le cadre de l'utilisation de buildroot, 
%n'est pas adapt\'ee. Buildroot ne fournit pas de solutions pour une telle manipulation et
%ne prends qu'un seul patch qui est cens\'e ajouter le support d'Ipipe. 
%Afin de disposer de Xenomai sur Red Pitaya, dans buildroot, pour la
%g\'en\'eration d'une image de mani\`ere automatique (sans interventions
%manuelles) la solution consiste donc \`a re-g\'en\'erer le patch Ipipe pour le
%noyau 4.9.24 avec le support du processeur cible.
%
%Ceci va se faire en 4 \'etapes~:
%% JMF TODO
%\begin{itemize}
%\item tel que d\'ecrit en section \ref{xe}, t\'el\'echarger un noyau propre et 
%les correctifs n\'ecessaires \`a la mise en \oe uvre de Ipipe.
%\item faire une copie des sources du noyau (non modifi\'ee) qui servira de r\'ef\'erence 
%pour \'etablir la diff\'erence entre version modifi\'ee et d'origine lors de l'application
%des patchs,
%\item se placer dans le r\'epertoire du noyau \`a modifier et appliquer les trois
%patchs tel que nous l'avons vu dans \ref{xe}
%\item r\'ealiser le diff entre la version non modifi\'ee et la version sur
%lesquels les patchs ont \'et\'es appliqu\'es \`a l'aide de la commande~:
%\begin{lstlisting}
%diff -ruN linux.orig linux > redpitaya-ipipe-core.patch
%\end{lstlisting}
%(avec 'r' pour r\'ecursif, 'u' pour
%unifi\'e (n\'ecessaire pour la commande patch) et 'N' pour que \textbf{diff}
%affiche le contenu complet du fichier s'il est absent d'un des deux r\'epertoires
%au lieu de simplement signaler qu'il n'est pr\'esent que dans l'un d'eux).
%\end{itemize}
%
%Cette s\'erie d'\'etapes a permis de disposer d'un fichier
%redpitaya-ipipe-core.patch \'equivalent \`a celui fourni par le projet 
%Xenomai mais avec le support pour Red Pitaya et applicable directement sur le noyau 
%s\'electionn\'e, condition n\'ecessaire au contexte buildroot.
%
%Ce fichier, par commodit\'e, peut \^etre d\'eplac\'e dans
%{\tt board/redpitaya/} d'une copie locale de buildroot tel que nous l'avons vu ici (dans
%notre cas, il s'agira de {\tt redpitaya/board/redpitaya}).
% 
%\subsection{Configuration de Buildroot et r\'ealisation d'un defconfig
%d\'edi\'e}
%
%Nous d\'esirons cr\'eer un fichier de configuration pour l'environnement buildroot 
%permettant la g\'en\'eration automatique d'une image avec le support temps-r\'eel.
%
%Le d\'epot fournit un support pour la Red Pitaya sans Xenomai.
%La premi\`ere \'etape consiste \`a configurer l'environnement par~:
%{\tt make redpitaya\_defconfig}.
%
%Ensuite, cette configuration par d\'efaut doit \^etre adapt\'ee aux contraintes
%li\'ees \`a la mise en place de Xenomai.
%
%Dans l'interface de configuration de buildroot les modifications sont les
%suivantes~:
%
%% JMF : TODO
%Parties noyau, sous menu \verb~Kernel --> :~, les op\'erations sont
%\begin{enumerate}
%\item
%s\'electionner le t\'el\'echargement d'un noyau sur d\'epot git~:
%\begin{itemize}
%\item Kernel Version est fix\'e \`a {\tt Custom Git repository} 
%%\^etre modifi\'e en {\tt custom tarball}
%%\item URL of custom kernel tarball doit \^etre renseign\'e avec la valeur : 
%%\begin{lstlisting}
%%https://github.com/linux-sunxi/linux-sunxi/archive/sunxi-v3.4.24-r2.tar.gz
%%\end{lstlisting}
%\end{itemize}
%\item
%le support pour l'extension Adeos/Ipipe doit \^etre activ\'ee en sp\'ecifiant le
%chemin relatif pour le patch g\'en\'er\'e pr\'ec\'edemment~:
%\begin{lstlisting}
%Kernel -->
%	Linux Kernel Extensions -->
%		[*] Adeos/Xenomai Real-Time patch
%		(board/redpitaya/redpitaya-ipipe-core.patch) Path for Adeos patch file)
%\end{lstlisting}
%\end{enumerate}
%
%Parties applicatives~:
%
%Xenomai et ses options de base sont activ\'es par d\'efaut (s\'electionn\'ees
%par l'activation de Adeos dans le menu kernel), seul testsuite doit \^etre
%activ\'ee manuellement~:
%\begin{lstlisting}
%Dans Target packages -->
%	Real-Time --> 
%		Xenomai Userspace
%		[*] Install testsuite
%		[*] Native skin library
%		[*] Posix skin library
%\end{lstlisting}
%
%Les modifications dans la configuration de buildroot \'etant finies, la
%derni\`ere \'etape consiste \`a produire le fichier de configuration par
%d\'efaut (defconfig) correspondant~:
%
%Apr\`es \^etre sortie de l'interface de configuration, la commande
%
%{\tt make savedefconfig} va produire un fichier
%defconfig. Ce fichier doit \^etre copi\'e en 
%{\tt configs/redpitaya\_xenomai\_defconfig}.

\newpage
\appendix
\section{Sending commands through a TCP/IP server}\label{tcp}

A basic example only allowing for a single connection (no multi-threading) of a TCP/IP 
server is as follows:

\lstinputlisting[basicstyle=\footnotesize\ttfamily,language=C]{./programs/tcp_server.c}

Reaching this server is for example achieved by running {\tt telnet IP 9999} for
connecting to port 9999 of the IP address of the server, or thanks to 
{\tt netcat} with {\tt echo "message | nc IP port}.

\section{Modifying the Buildroot configuration}

Standard tools needed for the lab sessions:

\begin{tabular}{l|p{0.46\linewidth}}\hline
User space& ({\tt make menuconfig}) \\\hline\hline
screen & (multiple terminals)\\ 
lighttpd & (web server) \\
sudo & (change user: to test www-data or execute as root CGI commands)\\
dropbear & (ssh server)\\
rsync & (synchronizing file systems)\\
file & (identifying file type)\\
stress & (loading the scheduler)\\
vim & (after activating {\em show packages that are also provided by busybox}) \\
nano & (basic editor) \\ \hline
\end{tabular}

Updates to the kernel configuration to allow accessing the GPIOs and LEDs, as well as
configuring the FPGA:

\begin{tabular}{l|p{0.70\linewidth}}\hline
Kernel space &({\tt make linux-menuconfig}) \\\hline\hline
LED driver as module & {\tt Device Drivers} $\rightarrow$ {\tt LED Support} $\rightarrow$ {\tt
$<$M$>$ LED Class Support} et {\tt $<$M$>$ LED Support for GPIO connected LED}\\ 
GPIO driver as module & {\tt Device Drivers} $\rightarrow$ {\tt GPIO Support} $\rightarrow$ 
{\tt Memory mapped GPIO drivers} $\rightarrow$ {\tt $<$M$>$ Xilinx GPIO support} et {\tt $<$M$>$ 
Xilinx Zynq GPIO support}\\
%xdevcfg as module & {\tt Device Drivers} $\rightarrow$ {\tt Character devices} $\rightarrow$ {\tt Xilinx Device Configuration}\\
%FPGA Manager as module & {\tt Device Drivers} $\rightarrow$ {\tt FPGA Configuration Support} 
%$\rightarrow$ {\tt 
%$<$M$>$ FPGA Configuration Framework} \\
\end{tabular}

\begin{thebibliography}{9}
\bibitem{ficheux} P. Ficheux, {\em Les distributions ``embarqu\'ees'' pour Rasberry PI}, 
Opensilicium 7 (2013)
\bibitem{xeno} P. Kadionik, P.Ficheux, {\em Temps r\'eel sous LINUX (reloaded)},
GNU/Linux Magazine France HS 24 (2006), disponible \`a 
\url{http://pficheux.free.fr/articles/lmf/hs24/realtime/linux_realtime_reloaded_final_www.pdf}
et \url{http://www.unixgarden.com/index.php/gnu-linux-magazine-hs/temps-reel-sous-linux-reloaded}
\bibitem{ldd3} J. Corbet, A. Rubini, \& G. Kroah-Hartman, {\em Linux Device Drivers, 3rd Ed.},
O'Reilly, disponible \`a \url{http://lwn.net/Kernel/LDD3/}
\bibitem{zynq} {\em Zynq-7000 All Programmable SoC Technical Reference Manual, rev. 6 Dec. 2017}, Xilinx (2017)
\bibitem{blaess} C. Blaess, {\em Interactions entre espace utilisateur, noyau et mat\'eriel}, GNU/Linux
Magazine France Hors S\'erie 87 (Nov.--D\'ec. 2016) 
\bibitem{blaes} C. Blaess, {\em Solutions temps r\'eel sous Linux}, Eyrolles (2012)
\end{thebibliography}

\appendix
\section*{Activating the X11 server}\label{VGA}

The Red Pitaya is not fitted with a graphical interface port. Nevertheless,
graphical applications can be executed, thanks to a X11 server running on the
Red Pitaya and displaying on the client running on the PC (X11 swaps the usual
convention of server and client naming convention).

Activating the graphical interface capability is achieved by running
in Buildroot the {\tt X.org} menu in {\tt Target Packages} $\rightarrow$
{\tt Graphic libraries and applications (graphic/text)}. Remember to install
{\tt xauth} in order to handle certificates allowing for exporting the graphical
interface from the Red Pitaya to the host PC.

\begin{figure}[h!tb]
\begin{center}
\includegraphics[width=0.45\linewidth]{xterm_redpi}
\includegraphics[width=0.45\linewidth]{dillo_redpi}
\end{center}
\caption{Left: Xterm executed on the Red Pitaya and displayed on the host PC.
Right: Dillo allows for displaying HTML web pages from the Red Pitaya.}
\label{console}
\end{figure}

\subsection{Qt5 graphical user Interface}

% 160109
% les exemples Qt se trouvent dans /usr/lib/qt/examples/
% /usr/lib/qt/examples/gui/analogclock/analogclock -platform linuxfb

Qt5 is a set of multiplatform C++ libraries compatible with multiple operating systems
(MS-Windows, X11 especially on GNU/Linux, OS X or iOS for Apple) and hardware platforms (x86, MIPS, 
ARM, SPARC ...). Despite is heavy burden on computational requirements as is the case for all graphical
interface, it emphasizes the potential benefits of running an operating system on the embedded
system.

Compiling Qt5 support on Buildroot requires C++ support, which is not active in the default configuration.
We make sure the whole toolchain is re-compiled with C++ support by deleting the {\tt output} ({\tt rm -rf output})
directory in the Buildroot archive and running again {\tt make}.

Furthermore, Qt5 fonts are not included by default, and requires calling
({\tt fontconfig}) whose computational requirements are too heavy for an embedded system.
The most efficient approach is to include static fonts by activating the
{\tt liberation} package in the {\tt Target packages} $\rightarrow$ {\tt
Fonts, cursors, icons, sounds and themes} menu and, on the Red Pitaya, to create on the
embedded platform the symbolic link between the location in which fonts are
stored, ({\tt /usr/share/fonts/liberation}), and the directory in which Qt5 expects
the fonts to be located ({\tt /usr/lib/fonts}):
\begin{verbatim}
cd /usr/lib
ln -s /usr/share/fonts/liberation/ fonts
\end{verbatim}

Once the system has been installed, we can develop and compile a basic test application.
The trivial program proposed (initially for Qt4, but also functional with Qt5) at
\url{http://doc.qt.io/qt-4.8/gettingstartedqt.html} is saved in a filename with a
{\tt .cpp} extension (Fig. \ref{texte}).

\begin{figure}[h!tb]
{\footnotesize
\begin{minipage}[t]{\linewidth}
\begin{minipage}{0.46\linewidth}
\lstinputlisting[language=C++]{./programs/1st_test.cpp}

\includegraphics[width=\linewidth]{1st_test.png}
\end{minipage}\hfill
\begin{minipage}{0.46\linewidth}
\lstinputlisting[language=C++]{./programs/2nd_test.cpp}

\includegraphics[width=\linewidth]{2nd_test.png}
\end{minipage}
\end{minipage}
}
\caption{Left: example of a text editor, one of the widgets provided by the Qt5 library.
Right: a slightly more advanced example with buttons.}
\label{texte}
\end{figure}

\noindent\fbox{\parbox{\linewidth}{
\includegraphics[width=.7cm]{danger} %\danger 
The extension of the filename is important: a {\tt .c} extension results in an attempt
to compile using {\tt gcc} which does not understand C++ syntax nor knows how to link with
the appropriate libraries.}}

Qt provides the means for generating dynamically the {\tt Makefile} from a configuration file whose 
extension is {\tt .pro}. This file is automatically generated by running {\tt qmake -project}. However,
we must make sure we run {\tt qmake} from the {\tt output/host/usr/bin} directory of Buildroot and
{\em not the qmake possibly installed on the host (PC).} Having generated the {\tt .pro} configuration
file, we add the Qt modules required by the proposed program, in our case the widgets module, with
\begin{verbatim}
QT += widgets
\end{verbatim}
to finally have a configuration file looking like
\begin{verbatim}
TEMPLATE = app
TARGET = Qt
INCLUDEPATH += .

QT += widgets

# Input
SOURCES += 2nd_test.cpp
\end{verbatim}

We run again {\tt \$BR/output/host/usr/bin/qmake} -- this time without the
{\tt -project} option -- to generate the Makefile, and finally the source code
is compiled to a binary executable with {\tt make}.

Once the compilation has been completed, two results are obtained:
\begin{itemize}
\item the GNU/Linux image (rootfs) for the ARM target including the Qt libraries.
\item the host (PC) now has the Qt tools for configuring and cross-compiling Qt5 
programs written in C++ and calling Qt libraries.
\end{itemize}

The program is executed on the Red Pitaya platform by specifying the graphical output,
here {\em xcb}: after transferring the binary executable to the Red Pitaya target
platform, we execute {\tt Qt -platform xcb} and, if all goes well, a window with a text editor
should be displayed (Fig. \ref{texte}, bottom)

% penser a nommer son fichier.cpp et non .c

Once the procedure for compiling a program has been understood, the Qt language and libraries
are documented with many examples given in the archive at {\tt 
\$BR/output/build/qt5base-5.*/examples}. We have hence checked that the examples
{\tt gui/analogclock} or {\tt widgets/digitalclock} are working properly (Fig. \ref{qtex}). 

\begin{figure}[h!tb]
\begin{center}
\includegraphics[height=5.8cm]{qt_analog.png}
\includegraphics[height=5.8cm]{qt_digital.png}
%\includegraphics[height=3.8cm]{20150905_181736.jpg}
\end{center}
\caption{Left: the {\tt analogclock} example. Right: the
{\tt digitalclock} example. Both examples are found in {\tt 
/usr/lib/qt/examples/}, and are executed on the {\tt xcb} platform specified
with the {\tt -plateform} argument when running the executable.}
\label{qtex}
\end{figure}

\end{document}

\newpage
\section*{Appendix: mounting NFS}



FPGA manager framework
fpga_manager fpga0: Xilinx Zynq FPGA Manager registered
fpga-region fpga-full: FPGA Region probed
random: crng init done
fpga_manager fpga0: fpga_mgr_unregister Xilinx Zynq FPGA Manager
xdevcfg f8007000.devcfg: ioremap 0xf8007000 to dea36000


redpitaya> ls -l /dev/xdevcfg 
crw-rw----    1 root     root      240,   0 Jan  1 13:47 /dev/xdevcfg


redpitaya> rmmod xilinx_devcfg
redpitaya> modprobe zynq_fpga
redpitaya> ls -l /sys/class/fpga_manager/fpga0/
--w-------    1 root     root          4096 Jan  1 13:52 firmware


cd tmp/ex_axi_gpio/ex_axi_gpio.runs/impl_1
vi system_wrapper.bif
bootgen -image boot_a53_bin_pl_all.bif -arch zynqmp -process_bitstream bin
bootgen -image system_wrapper.bif -arch zynq -process_bitstream bin

redpitaya> mv /tmp/system_wrapper.bit.bin /lib/firmware/system_wrapper.bit.bin 
redpitaya> echo "system_wrapper.bit.bin" > /sys/class/fpga_manager/fpga0/firmware 
dmesg
fpga_manager fpga0: writing system_wrapper.bit.bin to Xilinx Zynq FPGA Manager


redpitaya> cd iio\:device0/
redpitaya> cat in_temp0_raw
2441
redpitaya> cat in_temp0_raw
2674


redpitaya> devmem 0x41200004
0x000000FF                      # 1 = les broches sont en entree
redpitaya> devmem 0x41200000
0x00000000
redpitaya> devmem 0x41200004 8 0x00
redpitaya> devmem 0x41200000 8 0xff
redpitaya> devmem 0x41200000 8 0x00
redpitaya> devmem 0x41200000 8 0x55
redpitaya> devmem 0x41200000 8 0xAA
redpitaya> while true; do devmem 0x41200000 8 0xAA ; sleep 0.5;devmem 0x41200000 8 0x55;sleep 0.5; done



Exemple AXI GPIO :
------------------

Partant de 2PL :
- ajouter AXI GPIO, double cliquer pour passer de 32 `a 8 bits
- changer le nom du port en led_o
- fichier de contrainte ports.xdc
- generer HDL top
